%%%%%%%%%%%%%%%%%%%%%%%%%%%%%%%%%%%%%%%%%%%%%%%%%%%%%%%%%%%%
% 14.10 INTEGRAL TABLES
% The Composition of Advanced Mathematics
%%%%%%%%%%%%%%%%%%%%%%%%%%%%%%%%%%%%%%%%%%%%%%%%%%%%%%%%%%%%

\section{Integral Tables}
\label{sec:integral-tables}

\prerequisite{
Antiderivatives, derivative rules, algebraic manipulation,
trigonometric identities, exponential and logarithmic functions,
inverse trigonometric functions, hyperbolic functions,
limits, and basic substitution.
}

\learningobjectives{
By the end of this reference section, the reader should be able to:
\begin{itemize}
    \item recognize standard indefinite-integral notation;
    \item distinguish indefinite and definite integrals;
    \item use the constant-multiple and sum rules;
    \item recognize the reverse power rule;
    \item integrate exponential and logarithmic forms;
    \item integrate the six standard trigonometric derivative patterns;
    \item recognize inverse-trigonometric antiderivative forms;
    \item recognize hyperbolic antiderivative forms;
    \item use substitution;
    \item identify common logarithmic substitution patterns;
    \item evaluate standard trigonometric integrals;
    \item apply power-reduction identities when appropriate;
    \item integrate selected rational functions;
    \item use partial-fraction decompositions;
    \item recognize completing-the-square patterns;
    \item use integration by parts;
    \item recognize reduction formulas;
    \item evaluate common definite integrals;
    \item recognize improper integrals and convergence conditions;
    \item identify Gaussian, Gamma, and Beta integrals;
    \item use symmetry in definite integration;
    \item connect integrals with area, accumulation, probability,
          differential equations, Fourier analysis, and geometry;
    \item avoid common integration errors.
\end{itemize}
}

%%%%%%%%%%%%%%%%%%%%%%%%%%%%%%%%%%%%%%%%%%%%%%%%%%%%%%%%%%%%
\subsection{Integral Notation}
\label{subsec:integral-notation-reference}
%%%%%%%%%%%%%%%%%%%%%%%%%%%%%%%%%%%%%%%%%%%%%%%%%%%%%%%%%%%%

The indefinite integral

\[
\boxed{
\int f(x)\,dx
}
\]

represents the family of antiderivatives of \(f\).

If

\[
F'(x)=f(x),
\]

then

\[
\boxed{
\int f(x)\,dx
=
F(x)+C,
}
\]

where \(C\) is an arbitrary constant.

%%%%%%%%%%%%%%%%%%%%%%%%%%%%%%%%%%%%%%%%%%%%%%%%%%%%%%%%%%%%
\subsection{Definite Integral}
\label{subsec:definite-integral-reference}
%%%%%%%%%%%%%%%%%%%%%%%%%%%%%%%%%%%%%%%%%%%%%%%%%%%%%%%%%%%%

The definite integral from \(a\) to \(b\) is

\[
\boxed{
\int_a^b f(x)\,dx.
}
\]

Under suitable conditions, it represents signed accumulation over the
interval \([a,b]\).

%%%%%%%%%%%%%%%%%%%%%%%%%%%%%%%%%%%%%%%%%%%%%%%%%%%%%%%%%%%%
\subsection{Fundamental Theorem of Calculus}
\label{subsec:ftc-integral-reference}
%%%%%%%%%%%%%%%%%%%%%%%%%%%%%%%%%%%%%%%%%%%%%%%%%%%%%%%%%%%%

If \(f\) is continuous on \([a,b]\) and

\[
F'(x)=f(x),
\]

then

\[
\boxed{
\int_a^b f(x)\,dx
=
F(b)-F(a).
}
\]

This is the evaluation form of the Fundamental Theorem of Calculus.

%%%%%%%%%%%%%%%%%%%%%%%%%%%%%%%%%%%%%%%%%%%%%%%%%%%%%%%%%%%%
\subsection{Accumulation Form of the Fundamental Theorem}
\label{subsec:ftc-accumulation-reference}
%%%%%%%%%%%%%%%%%%%%%%%%%%%%%%%%%%%%%%%%%%%%%%%%%%%%%%%%%%%%

If

\[
A(x)
=
\int_a^x f(t)\,dt,
\]

then, under standard continuity assumptions,

\[
\boxed{
A'(x)=f(x).
}
\]

Thus differentiation and integration are inverse processes in this sense.

%%%%%%%%%%%%%%%%%%%%%%%%%%%%%%%%%%%%%%%%%%%%%%%%%%%%%%%%%%%%
\subsection{Constant Rule}
\label{subsec:constant-integral-rule}
%%%%%%%%%%%%%%%%%%%%%%%%%%%%%%%%%%%%%%%%%%%%%%%%%%%%%%%%%%%%

For constant \(c\),

\[
\boxed{
\int c\,dx
=
cx+C.
}
\]

%%%%%%%%%%%%%%%%%%%%%%%%%%%%%%%%%%%%%%%%%%%%%%%%%%%%%%%%%%%%
\subsection{Constant-Multiple Rule}
\label{subsec:constant-multiple-integral-rule}
%%%%%%%%%%%%%%%%%%%%%%%%%%%%%%%%%%%%%%%%%%%%%%%%%%%%%%%%%%%%

For constant \(c\),

\[
\boxed{
\int c f(x)\,dx
=
c\int f(x)\,dx.
}
\]

%%%%%%%%%%%%%%%%%%%%%%%%%%%%%%%%%%%%%%%%%%%%%%%%%%%%%%%%%%%%
\subsection{Sum Rule}
\label{subsec:sum-integral-rule}
%%%%%%%%%%%%%%%%%%%%%%%%%%%%%%%%%%%%%%%%%%%%%%%%%%%%%%%%%%%%

\[
\boxed{
\int
[f(x)+g(x)]
\,dx
=
\int f(x)\,dx
+
\int g(x)\,dx.
}
\]

%%%%%%%%%%%%%%%%%%%%%%%%%%%%%%%%%%%%%%%%%%%%%%%%%%%%%%%%%%%%
\subsection{Difference Rule}
\label{subsec:difference-integral-rule}
%%%%%%%%%%%%%%%%%%%%%%%%%%%%%%%%%%%%%%%%%%%%%%%%%%%%%%%%%%%%

\[
\boxed{
\int
[f(x)-g(x)]
\,dx
=
\int f(x)\,dx
-
\int g(x)\,dx.
}
\]

%%%%%%%%%%%%%%%%%%%%%%%%%%%%%%%%%%%%%%%%%%%%%%%%%%%%%%%%%%%%
\subsection{Power Rule}
\label{subsec:power-integral-rule}
%%%%%%%%%%%%%%%%%%%%%%%%%%%%%%%%%%%%%%%%%%%%%%%%%%%%%%%%%%%%

For

\[
n\neq-1,
\]

\[
\boxed{
\int x^n\,dx
=
\frac{x^{n+1}}{n+1}
+
C.
}
\]

More generally,

\[
\boxed{
\int
[u(x)]^n u'(x)\,dx
=
\frac{
[u(x)]^{n+1}
}{
n+1
}
+
C.
}
\]

%%%%%%%%%%%%%%%%%%%%%%%%%%%%%%%%%%%%%%%%%%%%%%%%%%%%%%%%%%%%
\subsection{The Exceptional Power \(n=-1\)}
\label{subsec:inverse-power-integral}
%%%%%%%%%%%%%%%%%%%%%%%%%%%%%%%%%%%%%%%%%%%%%%%%%%%%%%%%%%%%

When

\[
n=-1,
\]

the power-rule denominator would vanish.

Instead,

\[
\boxed{
\int\frac1x\,dx
=
\ln|x|+C.
}
\]

More generally,

\[
\boxed{
\int
\frac{
u'(x)
}{
u(x)
}
\,dx
=
\ln|u(x)|+C.
}
\]

%%%%%%%%%%%%%%%%%%%%%%%%%%%%%%%%%%%%%%%%%%%%%%%%%%%%%%%%%%%%
\subsection{Worked Example: Power Rule}
\label{subsec:worked-power-integral}
%%%%%%%%%%%%%%%%%%%%%%%%%%%%%%%%%%%%%%%%%%%%%%%%%%%%%%%%%%%%

\begin{workedexamplebox}{Integrate a Polynomial}

Evaluate

\[
\int
(3x^4-2x^2+5)
\,dx.
\]

Integrate term by term:

\[
3\int x^4\,dx
-
2\int x^2\,dx
+
5\int1\,dx.
\]

Therefore,

\begin{answerbox}
\[
\boxed{
\int
(3x^4-2x^2+5)
\,dx
=
\frac35x^5
-
\frac23x^3
+
5x
+
C.
}
\]
\end{answerbox}

\end{workedexamplebox}

%%%%%%%%%%%%%%%%%%%%%%%%%%%%%%%%%%%%%%%%%%%%%%%%%%%%%%%%%%%%
\subsection{Natural Exponential Integral}
\label{subsec:natural-exponential-integral}
%%%%%%%%%%%%%%%%%%%%%%%%%%%%%%%%%%%%%%%%%%%%%%%%%%%%%%%%%%%%

\[
\boxed{
\int e^x\,dx
=
e^x+C.
}
\]

More generally,

\[
\boxed{
\int
e^{u(x)}
u'(x)
\,dx
=
e^{u(x)}
+
C.
}
\]

%%%%%%%%%%%%%%%%%%%%%%%%%%%%%%%%%%%%%%%%%%%%%%%%%%%%%%%%%%%%
\subsection{Exponential With Constant Coefficient}
\label{subsec:exp-constant-coefficient-integral}
%%%%%%%%%%%%%%%%%%%%%%%%%%%%%%%%%%%%%%%%%%%%%%%%%%%%%%%%%%%%

For constant \(a\neq0\),

\[
\boxed{
\int e^{ax}\,dx
=
\frac1a e^{ax}+C.
}
\]

%%%%%%%%%%%%%%%%%%%%%%%%%%%%%%%%%%%%%%%%%%%%%%%%%%%%%%%%%%%%
\subsection{General Exponential Base}
\label{subsec:general-exponential-integral}
%%%%%%%%%%%%%%%%%%%%%%%%%%%%%%%%%%%%%%%%%%%%%%%%%%%%%%%%%%%%

For

\[
a>0,
\qquad
a\neq1,
\]

\[
\boxed{
\int a^x\,dx
=
\frac{
a^x
}{
\ln a
}
+
C.
}
\]

More generally,

\[
\boxed{
\int
a^{u(x)}
u'(x)\,dx
=
\frac{
a^{u(x)}
}{
\ln a
}
+
C.
}
\]

%%%%%%%%%%%%%%%%%%%%%%%%%%%%%%%%%%%%%%%%%%%%%%%%%%%%%%%%%%%%
\subsection{Integral of Natural Logarithm}
\label{subsec:integral-natural-logarithm}
%%%%%%%%%%%%%%%%%%%%%%%%%%%%%%%%%%%%%%%%%%%%%%%%%%%%%%%%%%%%

Using integration by parts,

\[
\boxed{
\int\ln x\,dx
=
x\ln x-x+C,
\qquad
x>0.
}
\]

More generally,

\[
\boxed{
\int\ln|x|\,dx
=
x\ln|x|-x+C,
\qquad
x\neq0.
}
\]

%%%%%%%%%%%%%%%%%%%%%%%%%%%%%%%%%%%%%%%%%%%%%%%%%%%%%%%%%%%%
\subsection{Sine Integral}
\label{subsec:sine-integral-reference}
%%%%%%%%%%%%%%%%%%%%%%%%%%%%%%%%%%%%%%%%%%%%%%%%%%%%%%%%%%%%

\[
\boxed{
\int\sin x\,dx
=
-\cos x+C.
}
\]

For constant \(a\neq0\),

\[
\boxed{
\int\sin(ax)\,dx
=
-\frac1a\cos(ax)+C.
}
\]

%%%%%%%%%%%%%%%%%%%%%%%%%%%%%%%%%%%%%%%%%%%%%%%%%%%%%%%%%%%%
\subsection{Cosine Integral}
\label{subsec:cosine-integral-reference}
%%%%%%%%%%%%%%%%%%%%%%%%%%%%%%%%%%%%%%%%%%%%%%%%%%%%%%%%%%%%

\[
\boxed{
\int\cos x\,dx
=
\sin x+C.
}
\]

For constant \(a\neq0\),

\[
\boxed{
\int\cos(ax)\,dx
=
\frac1a\sin(ax)+C.
}
\]

%%%%%%%%%%%%%%%%%%%%%%%%%%%%%%%%%%%%%%%%%%%%%%%%%%%%%%%%%%%%
\subsection{Secant-Squared Integral}
\label{subsec:secant-squared-integral}
%%%%%%%%%%%%%%%%%%%%%%%%%%%%%%%%%%%%%%%%%%%%%%%%%%%%%%%%%%%%

\[
\boxed{
\int\sec^2x\,dx
=
\tan x+C.
}
\]

More generally,

\[
\boxed{
\int
\sec^2(u)
u'\,dx
=
\tan u+C.
}
\]

%%%%%%%%%%%%%%%%%%%%%%%%%%%%%%%%%%%%%%%%%%%%%%%%%%%%%%%%%%%%
\subsection{Cosecant-Squared Integral}
\label{subsec:cosecant-squared-integral}
%%%%%%%%%%%%%%%%%%%%%%%%%%%%%%%%%%%%%%%%%%%%%%%%%%%%%%%%%%%%

\[
\boxed{
\int\csc^2x\,dx
=
-\cot x+C.
}
\]

%%%%%%%%%%%%%%%%%%%%%%%%%%%%%%%%%%%%%%%%%%%%%%%%%%%%%%%%%%%%
\subsection{Secant-Tangent Integral}
\label{subsec:secant-tangent-integral}
%%%%%%%%%%%%%%%%%%%%%%%%%%%%%%%%%%%%%%%%%%%%%%%%%%%%%%%%%%%%

\[
\boxed{
\int
\sec x\tan x
\,dx
=
\sec x+C.
}
\]

%%%%%%%%%%%%%%%%%%%%%%%%%%%%%%%%%%%%%%%%%%%%%%%%%%%%%%%%%%%%
\subsection{Cosecant-Cotangent Integral}
\label{subsec:cosecant-cotangent-integral}
%%%%%%%%%%%%%%%%%%%%%%%%%%%%%%%%%%%%%%%%%%%%%%%%%%%%%%%%%%%%

\[
\boxed{
\int
\csc x\cot x
\,dx
=
-\csc x+C.
}
\]

%%%%%%%%%%%%%%%%%%%%%%%%%%%%%%%%%%%%%%%%%%%%%%%%%%%%%%%%%%%%
\subsection{Tangent Integral}
\label{subsec:tangent-integral-reference}
%%%%%%%%%%%%%%%%%%%%%%%%%%%%%%%%%%%%%%%%%%%%%%%%%%%%%%%%%%%%

Since

\[
\tan x
=
\frac{\sin x}{\cos x},
\]

\[
\boxed{
\int\tan x\,dx
=
-\ln|\cos x|+C.
}
\]

Equivalently,

\[
\boxed{
\int\tan x\,dx
=
\ln|\sec x|+C.
}
\]

%%%%%%%%%%%%%%%%%%%%%%%%%%%%%%%%%%%%%%%%%%%%%%%%%%%%%%%%%%%%
\subsection{Cotangent Integral}
\label{subsec:cotangent-integral-reference}
%%%%%%%%%%%%%%%%%%%%%%%%%%%%%%%%%%%%%%%%%%%%%%%%%%%%%%%%%%%%

\[
\boxed{
\int\cot x\,dx
=
\ln|\sin x|+C.
}
\]

%%%%%%%%%%%%%%%%%%%%%%%%%%%%%%%%%%%%%%%%%%%%%%%%%%%%%%%%%%%%
\subsection{Secant Integral}
\label{subsec:secant-integral-reference}
%%%%%%%%%%%%%%%%%%%%%%%%%%%%%%%%%%%%%%%%%%%%%%%%%%%%%%%%%%%%

A standard antiderivative is

\[
\boxed{
\int\sec x\,dx
=
\ln|\sec x+\tan x|+C.
}
\]

%%%%%%%%%%%%%%%%%%%%%%%%%%%%%%%%%%%%%%%%%%%%%%%%%%%%%%%%%%%%
\subsection{Cosecant Integral}
\label{subsec:cosecant-integral-reference}
%%%%%%%%%%%%%%%%%%%%%%%%%%%%%%%%%%%%%%%%%%%%%%%%%%%%%%%%%%%%

A standard form is

\[
\boxed{
\int\csc x\,dx
=
\ln|\csc x-\cot x|+C.
}
\]

Equivalent logarithmic forms may also be used.

%%%%%%%%%%%%%%%%%%%%%%%%%%%%%%%%%%%%%%%%%%%%%%%%%%%%%%%%%%%%
\subsection{Quick Trigonometric Integral Table}
\label{subsec:quick-trig-integral-table}
%%%%%%%%%%%%%%%%%%%%%%%%%%%%%%%%%%%%%%%%%%%%%%%%%%%%%%%%%%%%

\[
\boxed{
\int\sin x\,dx
=
-\cos x+C
}
\]

\[
\boxed{
\int\cos x\,dx
=
\sin x+C
}
\]

\[
\boxed{
\int\sec^2x\,dx
=
\tan x+C
}
\]

\[
\boxed{
\int\csc^2x\,dx
=
-\cot x+C
}
\]

\[
\boxed{
\int\sec x\tan x\,dx
=
\sec x+C
}
\]

\[
\boxed{
\int\csc x\cot x\,dx
=
-\csc x+C
}
\]

\[
\boxed{
\int\tan x\,dx
=
-\ln|\cos x|+C
}
\]

\[
\boxed{
\int\cot x\,dx
=
\ln|\sin x|+C
}
\]

\[
\boxed{
\int\sec x\,dx
=
\ln|\sec x+\tan x|+C
}
\]

\[
\boxed{
\int\csc x\,dx
=
\ln|\csc x-\cot x|+C.
}
\]

%%%%%%%%%%%%%%%%%%%%%%%%%%%%%%%%%%%%%%%%%%%%%%%%%%%%%%%%%%%%
\subsection{Inverse-Sine Pattern}
\label{subsec:inverse-sine-integral-pattern}
%%%%%%%%%%%%%%%%%%%%%%%%%%%%%%%%%%%%%%%%%%%%%%%%%%%%%%%%%%%%

For \(a>0\),

\[
\boxed{
\int
\frac{
dx
}{
\sqrt{a^2-x^2}
}
=
\arcsin
\left(
\frac{x}{a}
\right)
+
C.
}
\]

%%%%%%%%%%%%%%%%%%%%%%%%%%%%%%%%%%%%%%%%%%%%%%%%%%%%%%%%%%%%
\subsection{Inverse-Tangent Pattern}
\label{subsec:inverse-tangent-integral-pattern}
%%%%%%%%%%%%%%%%%%%%%%%%%%%%%%%%%%%%%%%%%%%%%%%%%%%%%%%%%%%%

For \(a>0\),

\[
\boxed{
\int
\frac{
dx
}{
a^2+x^2
}
=
\frac1a
\arctan
\left(
\frac{x}{a}
\right)
+
C.
}
\]

%%%%%%%%%%%%%%%%%%%%%%%%%%%%%%%%%%%%%%%%%%%%%%%%%%%%%%%%%%%%
\subsection{Inverse-Secant Pattern}
\label{subsec:inverse-secant-integral-pattern}
%%%%%%%%%%%%%%%%%%%%%%%%%%%%%%%%%%%%%%%%%%%%%%%%%%%%%%%%%%%%

For \(a>0\),

\[
\boxed{
\int
\frac{
dx
}{
|x|\sqrt{x^2-a^2}
}
=
\frac1a
\operatorname{arcsec}
\left(
\frac{|x|}{a}
\right)
+
C
}
\]

on intervals where the relevant branch convention is fixed.

Equivalent inverse-trigonometric forms may be preferable depending on the
domain.

%%%%%%%%%%%%%%%%%%%%%%%%%%%%%%%%%%%%%%%%%%%%%%%%%%%%%%%%%%%%
\subsection{Alternative Inverse-Tangent Form}
\label{subsec:inverse-tangent-scaled-pattern}
%%%%%%%%%%%%%%%%%%%%%%%%%%%%%%%%%%%%%%%%%%%%%%%%%%%%%%%%%%%%

For positive \(a,b\),

\[
\boxed{
\int
\frac{
dx
}{
a^2+b^2x^2
}
=
\frac{
1
}{
ab
}
\arctan
\left(
\frac{bx}{a}
\right)
+
C.
}
\]

%%%%%%%%%%%%%%%%%%%%%%%%%%%%%%%%%%%%%%%%%%%%%%%%%%%%%%%%%%%%
\subsection{Inverse-Sine Composite Pattern}
\label{subsec:inverse-sine-composite-pattern}
%%%%%%%%%%%%%%%%%%%%%%%%%%%%%%%%%%%%%%%%%%%%%%%%%%%%%%%%%%%%

If \(u=u(x)\), then

\[
\boxed{
\int
\frac{
u'
}{
\sqrt{1-u^2}
}
\,dx
=
\arcsin u+C.
}
\]

%%%%%%%%%%%%%%%%%%%%%%%%%%%%%%%%%%%%%%%%%%%%%%%%%%%%%%%%%%%%
\subsection{Inverse-Tangent Composite Pattern}
\label{subsec:inverse-tangent-composite-pattern}
%%%%%%%%%%%%%%%%%%%%%%%%%%%%%%%%%%%%%%%%%%%%%%%%%%%%%%%%%%%%

\[
\boxed{
\int
\frac{
u'
}{
1+u^2
}
\,dx
=
\arctan u+C.
}
\]

%%%%%%%%%%%%%%%%%%%%%%%%%%%%%%%%%%%%%%%%%%%%%%%%%%%%%%%%%%%%
\subsection{Hyperbolic Sine Integral}
\label{subsec:sinh-integral-reference}
%%%%%%%%%%%%%%%%%%%%%%%%%%%%%%%%%%%%%%%%%%%%%%%%%%%%%%%%%%%%

\[
\boxed{
\int\sinh x\,dx
=
\cosh x+C.
}
\]

%%%%%%%%%%%%%%%%%%%%%%%%%%%%%%%%%%%%%%%%%%%%%%%%%%%%%%%%%%%%
\subsection{Hyperbolic Cosine Integral}
\label{subsec:cosh-integral-reference}
%%%%%%%%%%%%%%%%%%%%%%%%%%%%%%%%%%%%%%%%%%%%%%%%%%%%%%%%%%%%

\[
\boxed{
\int\cosh x\,dx
=
\sinh x+C.
}
\]

%%%%%%%%%%%%%%%%%%%%%%%%%%%%%%%%%%%%%%%%%%%%%%%%%%%%%%%%%%%%
\subsection{Hyperbolic Secant-Squared Integral}
\label{subsec:sech-squared-integral}
%%%%%%%%%%%%%%%%%%%%%%%%%%%%%%%%%%%%%%%%%%%%%%%%%%%%%%%%%%%%

\[
\boxed{
\int
\operatorname{sech}^2x
\,dx
=
\tanh x+C.
}
\]

%%%%%%%%%%%%%%%%%%%%%%%%%%%%%%%%%%%%%%%%%%%%%%%%%%%%%%%%%%%%
\subsection{Hyperbolic Cosecant-Squared Integral}
\label{subsec:csch-squared-integral}
%%%%%%%%%%%%%%%%%%%%%%%%%%%%%%%%%%%%%%%%%%%%%%%%%%%%%%%%%%%%

\[
\boxed{
\int
\operatorname{csch}^2x
\,dx
=
-\coth x+C.
}
\]

%%%%%%%%%%%%%%%%%%%%%%%%%%%%%%%%%%%%%%%%%%%%%%%%%%%%%%%%%%%%
\subsection{Hyperbolic Tangent Integral}
\label{subsec:tanh-integral-reference}
%%%%%%%%%%%%%%%%%%%%%%%%%%%%%%%%%%%%%%%%%%%%%%%%%%%%%%%%%%%%

\[
\boxed{
\int\tanh x\,dx
=
\ln(\cosh x)+C.
}
\]

%%%%%%%%%%%%%%%%%%%%%%%%%%%%%%%%%%%%%%%%%%%%%%%%%%%%%%%%%%%%
\subsection{Hyperbolic Cotangent Integral}
\label{subsec:coth-integral-reference}
%%%%%%%%%%%%%%%%%%%%%%%%%%%%%%%%%%%%%%%%%%%%%%%%%%%%%%%%%%%%

Where defined,

\[
\boxed{
\int\coth x\,dx
=
\ln|\sinh x|+C.
}
\]

%%%%%%%%%%%%%%%%%%%%%%%%%%%%%%%%%%%%%%%%%%%%%%%%%%%%%%%%%%%%
\subsection{Inverse Hyperbolic Sine Pattern}
\label{subsec:inverse-hyperbolic-sine-integral}
%%%%%%%%%%%%%%%%%%%%%%%%%%%%%%%%%%%%%%%%%%%%%%%%%%%%%%%%%%%%

For \(a>0\),

\[
\boxed{
\int
\frac{
dx
}{
\sqrt{x^2+a^2}
}
=
\operatorname{arsinh}
\left(
\frac{x}{a}
\right)
+
C.
}
\]

Equivalently,

\[
\boxed{
\int
\frac{
dx
}{
\sqrt{x^2+a^2}
}
=
\ln
\left|
x+\sqrt{x^2+a^2}
\right|
+
C.
}
\]

The two forms differ only by an additive constant depending on convention.

%%%%%%%%%%%%%%%%%%%%%%%%%%%%%%%%%%%%%%%%%%%%%%%%%%%%%%%%%%%%
\subsection{Inverse Hyperbolic Tangent Pattern}
\label{subsec:inverse-hyperbolic-tangent-integral}
%%%%%%%%%%%%%%%%%%%%%%%%%%%%%%%%%%%%%%%%%%%%%%%%%%%%%%%%%%%%

For \(|x|<a\),

\[
\boxed{
\int
\frac{
dx
}{
a^2-x^2
}
=
\frac{
1
}{
2a
}
\ln
\left|
\frac{
a+x
}{
a-x
}
\right|
+
C.
}
\]

Equivalently,

\[
\boxed{
\int
\frac{
dx
}{
a^2-x^2
}
=
\frac1a
\operatorname{artanh}
\left(
\frac{x}{a}
\right)
+
C
}
\]

on the appropriate interval.

%%%%%%%%%%%%%%%%%%%%%%%%%%%%%%%%%%%%%%%%%%%%%%%%%%%%%%%%%%%%
\subsection{Substitution Rule}
\label{subsec:substitution-integral-reference}
%%%%%%%%%%%%%%%%%%%%%%%%%%%%%%%%%%%%%%%%%%%%%%%%%%%%%%%%%%%%

If

\[
u=g(x),
\]

then

\[
du
=
g'(x)\,dx.
\]

Thus

\[
\boxed{
\int
f(g(x))
g'(x)
\,dx
=
\int
f(u)\,du.
}
\]

%%%%%%%%%%%%%%%%%%%%%%%%%%%%%%%%%%%%%%%%%%%%%%%%%%%%%%%%%%%%
\subsection{Worked Example: Basic Substitution}
\label{subsec:worked-basic-substitution}
%%%%%%%%%%%%%%%%%%%%%%%%%%%%%%%%%%%%%%%%%%%%%%%%%%%%%%%%%%%%

\begin{workedexamplebox}{Substitute an Inner Function}

Evaluate

\[
\int
2x(x^2+3)^5
\,dx.
\]

Let

\[
u=x^2+3.
\]

Then

\[
du=2x\,dx.
\]

Therefore,

\[
\int
2x(x^2+3)^5
\,dx
=
\int u^5\,du.
\]

Thus

\begin{answerbox}
\[
\boxed{
\int
2x(x^2+3)^5
\,dx
=
\frac{
(x^2+3)^6
}{
6
}
+
C.
}
\]
\end{answerbox}

\end{workedexamplebox}

%%%%%%%%%%%%%%%%%%%%%%%%%%%%%%%%%%%%%%%%%%%%%%%%%%%%%%%%%%%%
\subsection{Substitution Recognition Pattern}
\label{subsec:substitution-recognition-pattern}
%%%%%%%%%%%%%%%%%%%%%%%%%%%%%%%%%%%%%%%%%%%%%%%%%%%%%%%%%%%%

Look for expressions of form

\[
\boxed{
f(g(x))g'(x).
}
\]

Common signals include:

\[
\boxed{
u^n u',
}
\]

\[
\boxed{
e^u u',
}
\]

\[
\boxed{
\frac{u'}{u},
}
\]

\[
\boxed{
\cos u\,u',
}
\]

\[
\boxed{
\frac{u'}{1+u^2},
}
\]

and

\[
\boxed{
\frac{u'}{\sqrt{1-u^2}}.
}
\]

%%%%%%%%%%%%%%%%%%%%%%%%%%%%%%%%%%%%%%%%%%%%%%%%%%%%%%%%%%%%
\subsection{Definite Integral Substitution}
\label{subsec:definite-substitution-reference}
%%%%%%%%%%%%%%%%%%%%%%%%%%%%%%%%%%%%%%%%%%%%%%%%%%%%%%%%%%%%

If

\[
u=g(x),
\]

then one may either:

\begin{enumerate}
    \item change the bounds to \(u\)-values; or
    \item return to \(x\) before applying the original bounds.
\end{enumerate}

If \(x=a\) gives \(u=\alpha\) and \(x=b\) gives \(u=\beta\), then

\[
\boxed{
\int_a^b
f(g(x))g'(x)\,dx
=
\int_\alpha^\beta
f(u)\,du.
}
\]

%%%%%%%%%%%%%%%%%%%%%%%%%%%%%%%%%%%%%%%%%%%%%%%%%%%%%%%%%%%%
\subsection{Worked Example: Definite Substitution}
\label{subsec:worked-definite-substitution}
%%%%%%%%%%%%%%%%%%%%%%%%%%%%%%%%%%%%%%%%%%%%%%%%%%%%%%%%%%%%

\begin{workedexamplebox}{Change the Bounds}

Evaluate

\[
\int_0^1
2x e^{x^2}
\,dx.
\]

Let

\[
u=x^2.
\]

Then

\[
du=2x\,dx.
\]

The bounds become

\[
x=0
\Rightarrow
u=0,
\]

and

\[
x=1
\Rightarrow
u=1.
\]

Therefore,

\[
\int_0^1
2xe^{x^2}\,dx
=
\int_0^1
e^u\,du.
\]

Hence

\begin{answerbox}
\[
\boxed{
\int_0^1
2xe^{x^2}\,dx
=
e-1.
}
\]
\end{answerbox}

\end{workedexamplebox}

%%%%%%%%%%%%%%%%%%%%%%%%%%%%%%%%%%%%%%%%%%%%%%%%%%%%%%%%%%%%
\subsection{Integration by Parts}
\label{subsec:integration-by-parts-reference}
%%%%%%%%%%%%%%%%%%%%%%%%%%%%%%%%%%%%%%%%%%%%%%%%%%%%%%%%%%%%

From the product rule,

\[
(uv)'
=
u'v+uv',
\]

one obtains

\[
\boxed{
\int u\,dv
=
uv
-
\int v\,du.
}
\]

%%%%%%%%%%%%%%%%%%%%%%%%%%%%%%%%%%%%%%%%%%%%%%%%%%%%%%%%%%%%
\subsection{Choosing \(u\) in Integration by Parts}
\label{subsec:liate-integration-by-parts}
%%%%%%%%%%%%%%%%%%%%%%%%%%%%%%%%%%%%%%%%%%%%%%%%%%%%%%%%%%%%

A common heuristic for choosing \(u\) is the order

\[
\boxed{
\text{Logarithmic}
\rightarrow
\text{Inverse trig}
\rightarrow
\text{Algebraic}
\rightarrow
\text{Trigonometric}
\rightarrow
\text{Exponential}.
}
\]

This is often remembered as LIATE.

It is a heuristic, not a theorem.

%%%%%%%%%%%%%%%%%%%%%%%%%%%%%%%%%%%%%%%%%%%%%%%%%%%%%%%%%%%%
\subsection{Worked Example: Integration by Parts}
\label{subsec:worked-integration-by-parts}
%%%%%%%%%%%%%%%%%%%%%%%%%%%%%%%%%%%%%%%%%%%%%%%%%%%%%%%%%%%%

\begin{workedexamplebox}{Integrate \(xe^x\)}

Let

\[
u=x,
\qquad
dv=e^x\,dx.
\]

Then

\[
du=dx,
\qquad
v=e^x.
\]

Using integration by parts,

\[
\int xe^x\,dx
=
xe^x-\int e^x\,dx.
\]

Therefore,

\begin{answerbox}
\[
\boxed{
\int xe^x\,dx
=
xe^x-e^x+C
=
e^x(x-1)+C.
}
\]
\end{answerbox}

\end{workedexamplebox}

%%%%%%%%%%%%%%%%%%%%%%%%%%%%%%%%%%%%%%%%%%%%%%%%%%%%%%%%%%%%
\subsection{Integration by Parts for Logarithms}
\label{subsec:integration-by-parts-logarithm}
%%%%%%%%%%%%%%%%%%%%%%%%%%%%%%%%%%%%%%%%%%%%%%%%%%%%%%%%%%%%

To integrate

\[
\ln x,
\]

write

\[
\ln x
=
(\ln x)(1).
\]

Choose

\[
u=\ln x,
\qquad
dv=dx.
\]

Then

\[
du=\frac1x\,dx,
\qquad
v=x.
\]

Thus

\[
\boxed{
\int\ln x\,dx
=
x\ln x-x+C.
}
\]

%%%%%%%%%%%%%%%%%%%%%%%%%%%%%%%%%%%%%%%%%%%%%%%%%%%%%%%%%%%%
\subsection{Repeated Integration by Parts}
\label{subsec:repeated-integration-by-parts}
%%%%%%%%%%%%%%%%%%%%%%%%%%%%%%%%%%%%%%%%%%%%%%%%%%%%%%%%%%%%

Repeated integration by parts is useful for products such as

\[
x^n e^x,
\]

\[
x^n\sin x,
\]

and

\[
x^n\cos x.
\]

Each application lowers the polynomial degree.

%%%%%%%%%%%%%%%%%%%%%%%%%%%%%%%%%%%%%%%%%%%%%%%%%%%%%%%%%%%%
\subsection{Cyclic Integration by Parts}
\label{subsec:cyclic-integration-by-parts}
%%%%%%%%%%%%%%%%%%%%%%%%%%%%%%%%%%%%%%%%%%%%%%%%%%%%%%%%%%%%

Some integrals return to the original integral after repeated integration by
parts.

A standard example is

\[
\int e^x\sin x\,dx.
\]

Solving algebraically after two integrations by parts gives

\[
\boxed{
\int e^x\sin x\,dx
=
\frac12
e^x
(\sin x-\cos x)
+
C.
}
\]

Similarly,

\[
\boxed{
\int e^x\cos x\,dx
=
\frac12
e^x
(\sin x+\cos x)
+
C.
}
\]

%%%%%%%%%%%%%%%%%%%%%%%%%%%%%%%%%%%%%%%%%%%%%%%%%%%%%%%%%%%%
\subsection{General \(e^{ax}\sin bx\) Integral}
\label{subsec:general-exp-sine-integral}
%%%%%%%%%%%%%%%%%%%%%%%%%%%%%%%%%%%%%%%%%%%%%%%%%%%%%%%%%%%%

For constants \(a,b\), not both zero,

\[
\boxed{
\int
e^{ax}\sin bx
\,dx
=
\frac{
e^{ax}
}{
a^2+b^2
}
\left(
a\sin bx-b\cos bx
\right)
+
C.
}
\]

%%%%%%%%%%%%%%%%%%%%%%%%%%%%%%%%%%%%%%%%%%%%%%%%%%%%%%%%%%%%
\subsection{General \(e^{ax}\cos bx\) Integral}
\label{subsec:general-exp-cos-integral}
%%%%%%%%%%%%%%%%%%%%%%%%%%%%%%%%%%%%%%%%%%%%%%%%%%%%%%%%%%%%

\[
\boxed{
\int
e^{ax}\cos bx
\,dx
=
\frac{
e^{ax}
}{
a^2+b^2
}
\left(
a\cos bx+b\sin bx
\right)
+
C.
}
\]

%%%%%%%%%%%%%%%%%%%%%%%%%%%%%%%%%%%%%%%%%%%%%%%%%%%%%%%%%%%%
\subsection{Trigonometric Powers: Sine and Cosine}
\label{subsec:sine-cosine-powers-integrals}
%%%%%%%%%%%%%%%%%%%%%%%%%%%%%%%%%%%%%%%%%%%%%%%%%%%%%%%%%%%%

For integrals of form

\[
\int
\sin^m x
\cos^n x
\,dx,
\]

common strategies are:

\begin{itemize}
    \item if \(m\) is odd, save one factor of \(\sin x\);
    \item if \(n\) is odd, save one factor of \(\cos x\);
    \item if both powers are even, use power-reduction identities.
\end{itemize}

%%%%%%%%%%%%%%%%%%%%%%%%%%%%%%%%%%%%%%%%%%%%%%%%%%%%%%%%%%%%
\subsection{Power-Reduction Identities}
\label{subsec:power-reduction-integral-table}
%%%%%%%%%%%%%%%%%%%%%%%%%%%%%%%%%%%%%%%%%%%%%%%%%%%%%%%%%%%%

\[
\boxed{
\sin^2x
=
\frac{
1-\cos2x
}{2},
}
\]

and

\[
\boxed{
\cos^2x
=
\frac{
1+\cos2x
}{2}.
}
\]

These are especially useful when both sine and cosine powers are even.

%%%%%%%%%%%%%%%%%%%%%%%%%%%%%%%%%%%%%%%%%%%%%%%%%%%%%%%%%%%%
\subsection{Integral of \(\sin^2x\)}
\label{subsec:integral-sine-squared}
%%%%%%%%%%%%%%%%%%%%%%%%%%%%%%%%%%%%%%%%%%%%%%%%%%%%%%%%%%%%

Using power reduction,

\[
\sin^2x
=
\frac12
-
\frac12\cos2x.
\]

Therefore,

\[
\boxed{
\int\sin^2x\,dx
=
\frac{x}{2}
-
\frac{\sin2x}{4}
+
C.
}
\]

%%%%%%%%%%%%%%%%%%%%%%%%%%%%%%%%%%%%%%%%%%%%%%%%%%%%%%%%%%%%
\subsection{Integral of \(\cos^2x\)}
\label{subsec:integral-cosine-squared}
%%%%%%%%%%%%%%%%%%%%%%%%%%%%%%%%%%%%%%%%%%%%%%%%%%%%%%%%%%%%

\[
\boxed{
\int\cos^2x\,dx
=
\frac{x}{2}
+
\frac{\sin2x}{4}
+
C.
}
\]

%%%%%%%%%%%%%%%%%%%%%%%%%%%%%%%%%%%%%%%%%%%%%%%%%%%%%%%%%%%%
\subsection{Worked Example: Odd Sine Power}
\label{subsec:worked-odd-sine-power}
%%%%%%%%%%%%%%%%%%%%%%%%%%%%%%%%%%%%%%%%%%%%%%%%%%%%%%%%%%%%

\begin{workedexamplebox}{Integrate \(\sin^3x\cos^2x\)}

Write

\[
\sin^3x
=
\sin x(1-\cos^2x).
\]

Then

\[
\int
\sin^3x\cos^2x
\,dx
=
\int
\sin x
(1-\cos^2x)
\cos^2x
\,dx.
\]

Let

\[
u=\cos x,
\qquad
du=-\sin x\,dx.
\]

Thus

\[
-\int
(1-u^2)u^2
\,du.
\]

Therefore,

\begin{answerbox}
\[
\boxed{
\int
\sin^3x\cos^2x
\,dx
=
-\frac{u^3}{3}
+
\frac{u^5}{5}
+
C
}
\]

and hence

\[
\boxed{
\int
\sin^3x\cos^2x
\,dx
=
-\frac{\cos^3x}{3}
+
\frac{\cos^5x}{5}
+
C.
}
\]
\end{answerbox}

\end{workedexamplebox}

%%%%%%%%%%%%%%%%%%%%%%%%%%%%%%%%%%%%%%%%%%%%%%%%%%%%%%%%%%%%
\subsection{Tangent and Secant Powers}
\label{subsec:tangent-secant-power-integrals}
%%%%%%%%%%%%%%%%%%%%%%%%%%%%%%%%%%%%%%%%%%%%%%%%%%%%%%%%%%%%

For

\[
\int\tan^m x\sec^n x\,dx,
\]

common strategies include:

\begin{itemize}
    \item if \(n\) is even, save a factor \(\sec^2x\);
    \item if \(m\) is odd, save a factor \(\sec x\tan x\);
    \item use
    \[
    \tan^2x
    =
    \sec^2x-1.
    \]
\end{itemize}

%%%%%%%%%%%%%%%%%%%%%%%%%%%%%%%%%%%%%%%%%%%%%%%%%%%%%%%%%%%%
\subsection{Cotangent and Cosecant Powers}
\label{subsec:cotangent-cosecant-power-integrals}
%%%%%%%%%%%%%%%%%%%%%%%%%%%%%%%%%%%%%%%%%%%%%%%%%%%%%%%%%%%%

For

\[
\int\cot^m x\csc^n x\,dx,
\]

analogous identities are

\[
\boxed{
\cot^2x
=
\csc^2x-1
}
\]

and

\[
\boxed{
1+\cot^2x
=
\csc^2x.
}
\]

%%%%%%%%%%%%%%%%%%%%%%%%%%%%%%%%%%%%%%%%%%%%%%%%%%%%%%%%%%%%
\subsection{Product-to-Sum Integration}
\label{subsec:product-to-sum-integral-reference}
%%%%%%%%%%%%%%%%%%%%%%%%%%%%%%%%%%%%%%%%%%%%%%%%%%%%%%%%%%%%

Products of trigonometric functions with different frequencies can often be
simplified using product-to-sum identities.

For example,

\[
\boxed{
\sin A\cos B
=
\frac12
[
\sin(A+B)+\sin(A-B)
].
}
\]

%%%%%%%%%%%%%%%%%%%%%%%%%%%%%%%%%%%%%%%%%%%%%%%%%%%%%%%%%%%%
\subsection{Worked Example: Product-to-Sum Integration}
\label{subsec:worked-product-to-sum-integration}
%%%%%%%%%%%%%%%%%%%%%%%%%%%%%%%%%%%%%%%%%%%%%%%%%%%%%%%%%%%%

\begin{workedexamplebox}{Integrate \(\sin3x\cos2x\)}

Use

\[
\sin3x\cos2x
=
\frac12
[
\sin5x+\sin x
].
\]

Then

\[
\int
\sin3x\cos2x
\,dx
=
\frac12
\int
(\sin5x+\sin x)
\,dx.
\]

Therefore,

\begin{answerbox}
\[
\boxed{
\int
\sin3x\cos2x
\,dx
=
-\frac1{10}\cos5x
-
\frac12\cos x
+
C.
}
\]
\end{answerbox}

\end{workedexamplebox}

%%%%%%%%%%%%%%%%%%%%%%%%%%%%%%%%%%%%%%%%%%%%%%%%%%%%%%%%%%%%
\subsection{Trigonometric Substitution}
\label{subsec:trigonometric-substitution-reference}
%%%%%%%%%%%%%%%%%%%%%%%%%%%%%%%%%%%%%%%%%%%%%%%%%%%%%%%%%%%%

Common radical patterns suggest trigonometric substitutions.

For

\[
\boxed{
\sqrt{a^2-x^2},
}
\]

use

\[
\boxed{
x=a\sin\theta.
}
\]

For

\[
\boxed{
\sqrt{a^2+x^2},
}
\]

use

\[
\boxed{
x=a\tan\theta.
}
\]

For

\[
\boxed{
\sqrt{x^2-a^2},
}
\]

use

\[
\boxed{
x=a\sec\theta.
}
\]

%%%%%%%%%%%%%%%%%%%%%%%%%%%%%%%%%%%%%%%%%%%%%%%%%%%%%%%%%%%%
\subsection{Why the Standard Trigonometric Substitutions Work}
\label{subsec:why-trig-substitution-works}
%%%%%%%%%%%%%%%%%%%%%%%%%%%%%%%%%%%%%%%%%%%%%%%%%%%%%%%%%%%%

They exploit the Pythagorean identities:

\[
\boxed{
1-\sin^2\theta
=
\cos^2\theta,
}
\]

\[
\boxed{
1+\tan^2\theta
=
\sec^2\theta,
}
\]

and

\[
\boxed{
\sec^2\theta-1
=
\tan^2\theta.
}
\]

%%%%%%%%%%%%%%%%%%%%%%%%%%%%%%%%%%%%%%%%%%%%%%%%%%%%%%%%%%%%
\subsection{Worked Example: Trigonometric Substitution}
\label{subsec:worked-trig-substitution}
%%%%%%%%%%%%%%%%%%%%%%%%%%%%%%%%%%%%%%%%%%%%%%%%%%%%%%%%%%%%

\begin{workedexamplebox}{Integrate a Circular Radical}

Evaluate

\[
\int
\sqrt{a^2-x^2}
\,dx,
\qquad
a>0.
\]

Let

\[
x=a\sin\theta.
\]

Then

\[
dx=a\cos\theta\,d\theta,
\]

and

\[
\sqrt{a^2-x^2}
=
a\cos\theta
\]

on the chosen interval.

Thus

\[
\int
\sqrt{a^2-x^2}
\,dx
=
a^2
\int
\cos^2\theta
\,d\theta.
\]

Using power reduction and converting back,

\begin{answerbox}
\[
\boxed{
\int
\sqrt{a^2-x^2}
\,dx
=
\frac{x}{2}
\sqrt{a^2-x^2}
+
\frac{a^2}{2}
\arcsin
\left(
\frac{x}{a}
\right)
+
C.
}
\]
\end{answerbox}

\end{workedexamplebox}

%%%%%%%%%%%%%%%%%%%%%%%%%%%%%%%%%%%%%%%%%%%%%%%%%%%%%%%%%%%%
\subsection{Completing the Square Before Integration}
\label{subsec:complete-square-integral}
%%%%%%%%%%%%%%%%%%%%%%%%%%%%%%%%%%%%%%%%%%%%%%%%%%%%%%%%%%%%

Quadratic denominators often become standard forms after completing the
square.

For example,

\[
x^2+4x+13
\]

becomes

\[
\boxed{
(x+2)^2+9.
}
\]

Thus

\[
\int
\frac{
dx
}{
x^2+4x+13
}
\]

can be converted to an inverse-tangent form.

%%%%%%%%%%%%%%%%%%%%%%%%%%%%%%%%%%%%%%%%%%%%%%%%%%%%%%%%%%%%
\subsection{Worked Example: Completing the Square}
\label{subsec:worked-complete-square-integral}
%%%%%%%%%%%%%%%%%%%%%%%%%%%%%%%%%%%%%%%%%%%%%%%%%%%%%%%%%%%%

\begin{workedexamplebox}{Quadratic Denominator}

Evaluate

\[
\int
\frac{
dx
}{
x^2+4x+13
}.
\]

Complete the square:

\[
x^2+4x+13
=
(x+2)^2+9.
\]

Let

\[
u=x+2.
\]

Then

\[
du=dx.
\]

Therefore,

\[
\int
\frac{
du
}{
u^2+3^2
}
=
\frac13
\arctan
\left(
\frac{u}{3}
\right)
+
C.
\]

Hence

\begin{answerbox}
\[
\boxed{
\int
\frac{
dx
}{
x^2+4x+13
}
=
\frac13
\arctan
\left(
\frac{x+2}{3}
\right)
+
C.
}
\]
\end{answerbox}

\end{workedexamplebox}

%%%%%%%%%%%%%%%%%%%%%%%%%%%%%%%%%%%%%%%%%%%%%%%%%%%%%%%%%%%%
\subsection{Partial Fractions}
\label{subsec:partial-fractions-integral-reference}
%%%%%%%%%%%%%%%%%%%%%%%%%%%%%%%%%%%%%%%%%%%%%%%%%%%%%%%%%%%%

Rational functions

\[
\frac{P(x)}{Q(x)}
\]

with

\[
\deg P<\deg Q
\]

can often be decomposed into simpler fractions when \(Q\) factors.

If

\[
\deg P\geq\deg Q,
\]

perform polynomial division first.

%%%%%%%%%%%%%%%%%%%%%%%%%%%%%%%%%%%%%%%%%%%%%%%%%%%%%%%%%%%%
\subsection{Distinct Linear Factors}
\label{subsec:partial-fractions-distinct-linear}
%%%%%%%%%%%%%%%%%%%%%%%%%%%%%%%%%%%%%%%%%%%%%%%%%%%%%%%%%%%%

For

\[
Q(x)
=
(x-a)(x-b),
\qquad
a\neq b,
\]

use

\[
\boxed{
\frac{P(x)}{
(x-a)(x-b)
}
=
\frac{A}{x-a}
+
\frac{B}{x-b}.
}
\]

Each term integrates logarithmically.

%%%%%%%%%%%%%%%%%%%%%%%%%%%%%%%%%%%%%%%%%%%%%%%%%%%%%%%%%%%%
\subsection{Repeated Linear Factors}
\label{subsec:partial-fractions-repeated-linear}
%%%%%%%%%%%%%%%%%%%%%%%%%%%%%%%%%%%%%%%%%%%%%%%%%%%%%%%%%%%%

For factor

\[
(x-a)^m,
\]

include

\[
\boxed{
\frac{A_1}{x-a}
+
\frac{A_2}{(x-a)^2}
+
\cdots
+
\frac{A_m}{(x-a)^m}.
}
\]

%%%%%%%%%%%%%%%%%%%%%%%%%%%%%%%%%%%%%%%%%%%%%%%%%%%%%%%%%%%%
\subsection{Irreducible Quadratic Factors}
\label{subsec:partial-fractions-irreducible-quadratic}
%%%%%%%%%%%%%%%%%%%%%%%%%%%%%%%%%%%%%%%%%%%%%%%%%%%%%%%%%%%%

For an irreducible quadratic

\[
ax^2+bx+c,
\]

use a linear numerator:

\[
\boxed{
\frac{
Ax+B
}{
ax^2+bx+c
}.
}
\]

Repeated irreducible quadratics require corresponding repeated terms.

%%%%%%%%%%%%%%%%%%%%%%%%%%%%%%%%%%%%%%%%%%%%%%%%%%%%%%%%%%%%
\subsection{Worked Example: Partial Fractions}
\label{subsec:worked-partial-fractions}
%%%%%%%%%%%%%%%%%%%%%%%%%%%%%%%%%%%%%%%%%%%%%%%%%%%%%%%%%%%%

\begin{workedexamplebox}{Integrate a Rational Function}

Evaluate

\[
\int
\frac{
1
}{
x^2-1
}
\,dx.
\]

Factor:

\[
x^2-1
=
(x-1)(x+1).
\]

Decompose:

\[
\frac{
1
}{
x^2-1
}
=
\frac12
\frac1{x-1}
-
\frac12
\frac1{x+1}.
\]

Therefore,

\begin{answerbox}
\[
\boxed{
\int
\frac{
1
}{
x^2-1
}
\,dx
=
\frac12
\ln|x-1|
-
\frac12
\ln|x+1|
+
C.
}
\]
\end{answerbox}

\end{workedexamplebox}

%%%%%%%%%%%%%%%%%%%%%%%%%%%%%%%%%%%%%%%%%%%%%%%%%%%%%%%%%%%%
\subsection{Useful Quadratic Rational Pattern}
\label{subsec:quadratic-rational-pattern}
%%%%%%%%%%%%%%%%%%%%%%%%%%%%%%%%%%%%%%%%%%%%%%%%%%%%%%%%%%%%

When the numerator resembles the derivative of the denominator,

\[
\boxed{
\int
\frac{
2x+b
}{
x^2+bx+c
}
\,dx
=
\ln
|x^2+bx+c|
+
C
}
\]

where the denominator is nonzero.

%%%%%%%%%%%%%%%%%%%%%%%%%%%%%%%%%%%%%%%%%%%%%%%%%%%%%%%%%%%%
\subsection{Splitting a Linear Numerator}
\label{subsec:splitting-linear-numerator}
%%%%%%%%%%%%%%%%%%%%%%%%%%%%%%%%%%%%%%%%%%%%%%%%%%%%%%%%%%%%

For

\[
\int
\frac{
mx+n
}{
x^2+bx+c
}
\,dx,
\]

write

\[
mx+n
=
A(2x+b)+B.
\]

The integral then separates into:

\[
\boxed{
\text{a logarithmic term}
}
\]

and a remaining standard quadratic-denominator term.

%%%%%%%%%%%%%%%%%%%%%%%%%%%%%%%%%%%%%%%%%%%%%%%%%%%%%%%%%%%%
\subsection{Weierstrass Substitution}
\label{subsec:weierstrass-integral-reference}
%%%%%%%%%%%%%%%%%%%%%%%%%%%%%%%%%%%%%%%%%%%%%%%%%%%%%%%%%%%%

For rational functions of

\[
\sin x
\]

and

\[
\cos x,
\]

one may use

\[
\boxed{
t
=
\tan\left(\frac{x}{2}\right).
}
\]

Then

\[
\boxed{
\sin x
=
\frac{
2t
}{
1+t^2
},
}
\]

\[
\boxed{
\cos x
=
\frac{
1-t^2
}{
1+t^2
},
}
\]

and

\[
\boxed{
dx
=
\frac{
2
}{
1+t^2
}
\,dt.
}
\]

This converts a rational trigonometric integral into a rational function of
\(t\).

%%%%%%%%%%%%%%%%%%%%%%%%%%%%%%%%%%%%%%%%%%%%%%%%%%%%%%%%%%%%
\subsection{Reduction Formula for Powers of Sine}
\label{subsec:sine-reduction-formula}
%%%%%%%%%%%%%%%%%%%%%%%%%%%%%%%%%%%%%%%%%%%%%%%%%%%%%%%%%%%%

For

\[
n>1,
\]

a standard reduction formula is

\[
\boxed{
\int\sin^n x\,dx
=
-\frac{
\sin^{n-1}x\cos x
}{
n
}
+
\frac{
n-1
}{
n
}
\int
\sin^{n-2}x
\,dx.
}
\]

%%%%%%%%%%%%%%%%%%%%%%%%%%%%%%%%%%%%%%%%%%%%%%%%%%%%%%%%%%%%
\subsection{Reduction Formula for Powers of Cosine}
\label{subsec:cosine-reduction-formula}
%%%%%%%%%%%%%%%%%%%%%%%%%%%%%%%%%%%%%%%%%%%%%%%%%%%%%%%%%%%%

For

\[
n>1,
\]

\[
\boxed{
\int\cos^n x\,dx
=
\frac{
\cos^{n-1}x\sin x
}{
n
}
+
\frac{
n-1
}{
n
}
\int
\cos^{n-2}x
\,dx.
}
\]

%%%%%%%%%%%%%%%%%%%%%%%%%%%%%%%%%%%%%%%%%%%%%%%%%%%%%%%%%%%%
\subsection{Reduction Formula for \(x^n e^{ax}\)}
\label{subsec:power-exp-reduction}
%%%%%%%%%%%%%%%%%%%%%%%%%%%%%%%%%%%%%%%%%%%%%%%%%%%%%%%%%%%%

For \(a\neq0\),

\[
\boxed{
\int
x^n e^{ax}
\,dx
=
\frac{
x^n e^{ax}
}{
a
}
-
\frac{
n
}{
a
}
\int
x^{n-1}e^{ax}
\,dx.
}
\]

%%%%%%%%%%%%%%%%%%%%%%%%%%%%%%%%%%%%%%%%%%%%%%%%%%%%%%%%%%%%
\subsection{Definite Integral Reversal}
\label{subsec:definite-integral-reversal}
%%%%%%%%%%%%%%%%%%%%%%%%%%%%%%%%%%%%%%%%%%%%%%%%%%%%%%%%%%%%

\[
\boxed{
\int_a^b f(x)\,dx
=
-
\int_b^a f(x)\,dx.
}
\]

%%%%%%%%%%%%%%%%%%%%%%%%%%%%%%%%%%%%%%%%%%%%%%%%%%%%%%%%%%%%
\subsection{Zero-Length Interval}
\label{subsec:zero-length-integral}
%%%%%%%%%%%%%%%%%%%%%%%%%%%%%%%%%%%%%%%%%%%%%%%%%%%%%%%%%%%%

\[
\boxed{
\int_a^a f(x)\,dx
=
0.
}
\]

%%%%%%%%%%%%%%%%%%%%%%%%%%%%%%%%%%%%%%%%%%%%%%%%%%%%%%%%%%%%
\subsection{Interval Additivity}
\label{subsec:definite-integral-additivity}
%%%%%%%%%%%%%%%%%%%%%%%%%%%%%%%%%%%%%%%%%%%%%%%%%%%%%%%%%%%%

For any intermediate point \(c\),

\[
\boxed{
\int_a^b f(x)\,dx
=
\int_a^c f(x)\,dx
+
\int_c^b f(x)\,dx.
}
\]

%%%%%%%%%%%%%%%%%%%%%%%%%%%%%%%%%%%%%%%%%%%%%%%%%%%%%%%%%%%%
\subsection{Linearity of Definite Integrals}
\label{subsec:definite-integral-linearity}
%%%%%%%%%%%%%%%%%%%%%%%%%%%%%%%%%%%%%%%%%%%%%%%%%%%%%%%%%%%%

For constants \(A,B\),

\[
\boxed{
\int_a^b
[Af(x)+Bg(x)]
\,dx
=
A\int_a^b f(x)\,dx
+
B\int_a^b g(x)\,dx.
}
\]

%%%%%%%%%%%%%%%%%%%%%%%%%%%%%%%%%%%%%%%%%%%%%%%%%%%%%%%%%%%%
\subsection{Comparison Property}
\label{subsec:definite-integral-comparison}
%%%%%%%%%%%%%%%%%%%%%%%%%%%%%%%%%%%%%%%%%%%%%%%%%%%%%%%%%%%%

If

\[
f(x)\leq g(x)
\]

for every \(x\in[a,b]\), then

\[
\boxed{
\int_a^b f(x)\,dx
\leq
\int_a^b g(x)\,dx.
}
\]

%%%%%%%%%%%%%%%%%%%%%%%%%%%%%%%%%%%%%%%%%%%%%%%%%%%%%%%%%%%%
\subsection{Absolute-Value Bound}
\label{subsec:integral-absolute-value-bound}
%%%%%%%%%%%%%%%%%%%%%%%%%%%%%%%%%%%%%%%%%%%%%%%%%%%%%%%%%%%%

For integrable \(f\),

\[
\boxed{
\left|
\int_a^b f(x)\,dx
\right|
\leq
\int_a^b |f(x)|\,dx.
}
\]

%%%%%%%%%%%%%%%%%%%%%%%%%%%%%%%%%%%%%%%%%%%%%%%%%%%%%%%%%%%%
\subsection{Even Functions on Symmetric Intervals}
\label{subsec:even-function-integral-symmetry}
%%%%%%%%%%%%%%%%%%%%%%%%%%%%%%%%%%%%%%%%%%%%%%%%%%%%%%%%%%%%

If \(f\) is even,

\[
f(-x)=f(x),
\]

then

\[
\boxed{
\int_{-a}^{a}
f(x)\,dx
=
2
\int_0^a
f(x)\,dx.
}
\]

%%%%%%%%%%%%%%%%%%%%%%%%%%%%%%%%%%%%%%%%%%%%%%%%%%%%%%%%%%%%
\subsection{Odd Functions on Symmetric Intervals}
\label{subsec:odd-function-integral-symmetry}
%%%%%%%%%%%%%%%%%%%%%%%%%%%%%%%%%%%%%%%%%%%%%%%%%%%%%%%%%%%%

If \(f\) is odd,

\[
f(-x)=-f(x),
\]

then

\[
\boxed{
\int_{-a}^{a}
f(x)\,dx
=
0.
}
\]

%%%%%%%%%%%%%%%%%%%%%%%%%%%%%%%%%%%%%%%%%%%%%%%%%%%%%%%%%%%%
\subsection{Worked Example: Symmetry}
\label{subsec:worked-integral-symmetry}
%%%%%%%%%%%%%%%%%%%%%%%%%%%%%%%%%%%%%%%%%%%%%%%%%%%%%%%%%%%%

\begin{workedexamplebox}{Use Odd Symmetry}

Evaluate

\[
\int_{-2}^{2}
x^3\cos x
\,dx.
\]

The function

\[
x^3
\]

is odd, while

\[
\cos x
\]

is even.

Thus

\[
x^3\cos x
\]

is odd.

Therefore,

\begin{answerbox}
\[
\boxed{
\int_{-2}^{2}
x^3\cos x
\,dx
=
0.
}
\]
\end{answerbox}

\end{workedexamplebox}

%%%%%%%%%%%%%%%%%%%%%%%%%%%%%%%%%%%%%%%%%%%%%%%%%%%%%%%%%%%%
\subsection{Periodic Function Integral}
\label{subsec:periodic-integral-reference}
%%%%%%%%%%%%%%%%%%%%%%%%%%%%%%%%%%%%%%%%%%%%%%%%%%%%%%%%%%%%

If \(f\) has period \(T\), then

\[
\boxed{
\int_a^{a+T}
f(x)\,dx
}
\]

has the same value for every \(a\), assuming integrability.

Thus

\[
\boxed{
\int_a^{a+T}
f(x)\,dx
=
\int_0^T
f(x)\,dx.
}
\]

%%%%%%%%%%%%%%%%%%%%%%%%%%%%%%%%%%%%%%%%%%%%%%%%%%%%%%%%%%%%
\subsection{Mean Value of a Function}
\label{subsec:average-value-integral-reference}
%%%%%%%%%%%%%%%%%%%%%%%%%%%%%%%%%%%%%%%%%%%%%%%%%%%%%%%%%%%%

The average value of \(f\) on \([a,b]\) is

\[
\boxed{
f_{\mathrm{avg}}
=
\frac{
1
}{
b-a
}
\int_a^b
f(x)\,dx.
}
\]

%%%%%%%%%%%%%%%%%%%%%%%%%%%%%%%%%%%%%%%%%%%%%%%%%%%%%%%%%%%%
\subsection{Mean Value Theorem for Integrals}
\label{subsec:mean-value-theorem-integral-reference}
%%%%%%%%%%%%%%%%%%%%%%%%%%%%%%%%%%%%%%%%%%%%%%%%%%%%%%%%%%%%

If \(f\) is continuous on \([a,b]\), then there exists

\[
c\in[a,b]
\]

such that

\[
\boxed{
\int_a^b f(x)\,dx
=
f(c)(b-a).
}
\]

%%%%%%%%%%%%%%%%%%%%%%%%%%%%%%%%%%%%%%%%%%%%%%%%%%%%%%%%%%%%
\subsection{Area Between Curves}
\label{subsec:area-between-curves-integral}
%%%%%%%%%%%%%%%%%%%%%%%%%%%%%%%%%%%%%%%%%%%%%%%%%%%%%%%%%%%%

If

\[
f(x)\geq g(x)
\]

on \([a,b]\), then the area between the curves is

\[
\boxed{
A
=
\int_a^b
[f(x)-g(x)]
\,dx.
}
\]

For curves expressed in terms of \(y\),

\[
\boxed{
A
=
\int_c^d
[x_{\mathrm{right}}(y)-x_{\mathrm{left}}(y)]
\,dy.
}
\]

%%%%%%%%%%%%%%%%%%%%%%%%%%%%%%%%%%%%%%%%%%%%%%%%%%%%%%%%%%%%
\subsection{Volume by Disks}
\label{subsec:disk-method-integral-reference}
%%%%%%%%%%%%%%%%%%%%%%%%%%%%%%%%%%%%%%%%%%%%%%%%%%%%%%%%%%%%

For disks perpendicular to an axis,

\[
\boxed{
V
=
\pi
\int_a^b
R(x)^2
\,dx.
}
\]

%%%%%%%%%%%%%%%%%%%%%%%%%%%%%%%%%%%%%%%%%%%%%%%%%%%%%%%%%%%%
\subsection{Volume by Washers}
\label{subsec:washer-method-integral-reference}
%%%%%%%%%%%%%%%%%%%%%%%%%%%%%%%%%%%%%%%%%%%%%%%%%%%%%%%%%%%%

With outer radius \(R(x)\) and inner radius \(r(x)\),

\[
\boxed{
V
=
\pi
\int_a^b
\left[
R(x)^2-r(x)^2
\right]
dx.
}
\]

%%%%%%%%%%%%%%%%%%%%%%%%%%%%%%%%%%%%%%%%%%%%%%%%%%%%%%%%%%%%
\subsection{Volume by Cylindrical Shells}
\label{subsec:shell-method-integral-reference}
%%%%%%%%%%%%%%%%%%%%%%%%%%%%%%%%%%%%%%%%%%%%%%%%%%%%%%%%%%%%

For shell radius \(r(x)\) and height \(h(x)\),

\[
\boxed{
V
=
2\pi
\int_a^b
r(x)h(x)
\,dx.
}
\]

%%%%%%%%%%%%%%%%%%%%%%%%%%%%%%%%%%%%%%%%%%%%%%%%%%%%%%%%%%%%
\subsection{Arc Length}
\label{subsec:arc-length-integral-reference}
%%%%%%%%%%%%%%%%%%%%%%%%%%%%%%%%%%%%%%%%%%%%%%%%%%%%%%%%%%%%

For

\[
y=f(x),
\]

the arc length from \(a\) to \(b\) is

\[
\boxed{
L
=
\int_a^b
\sqrt{
1+[f'(x)]^2
}
\,dx.
}
\]

%%%%%%%%%%%%%%%%%%%%%%%%%%%%%%%%%%%%%%%%%%%%%%%%%%%%%%%%%%%%
\subsection{Parametric Arc Length}
\label{subsec:parametric-arc-length-integral}
%%%%%%%%%%%%%%%%%%%%%%%%%%%%%%%%%%%%%%%%%%%%%%%%%%%%%%%%%%%%

For

\[
x=x(t),
\qquad
y=y(t),
\]

\[
\boxed{
L
=
\int_\alpha^\beta
\sqrt{
\left(
\frac{dx}{dt}
\right)^2
+
\left(
\frac{dy}{dt}
\right)^2
}
\,dt.
}
\]

%%%%%%%%%%%%%%%%%%%%%%%%%%%%%%%%%%%%%%%%%%%%%%%%%%%%%%%%%%%%
\subsection{Polar Area}
\label{subsec:polar-area-integral-reference}
%%%%%%%%%%%%%%%%%%%%%%%%%%%%%%%%%%%%%%%%%%%%%%%%%%%%%%%%%%%%

For

\[
r=f(\theta),
\]

the enclosed area over \([\alpha,\beta]\) is

\[
\boxed{
A
=
\frac12
\int_\alpha^\beta
r^2\,d\theta.
}
\]

%%%%%%%%%%%%%%%%%%%%%%%%%%%%%%%%%%%%%%%%%%%%%%%%%%%%%%%%%%%%
\subsection{Improper Integrals}
\label{subsec:improper-integrals-reference}
%%%%%%%%%%%%%%%%%%%%%%%%%%%%%%%%%%%%%%%%%%%%%%%%%%%%%%%%%%%%

An integral is improper if:

\begin{itemize}
    \item at least one interval endpoint is infinite; or
    \item the integrand becomes unbounded on the interval.
\end{itemize}

Improper integrals are defined using limits.

%%%%%%%%%%%%%%%%%%%%%%%%%%%%%%%%%%%%%%%%%%%%%%%%%%%%%%%%%%%%
\subsection{Infinite-Interval Integral}
\label{subsec:infinite-interval-integral}
%%%%%%%%%%%%%%%%%%%%%%%%%%%%%%%%%%%%%%%%%%%%%%%%%%%%%%%%%%%%

For example,

\[
\boxed{
\int_a^\infty
f(x)\,dx
=
\lim_{b\to\infty}
\int_a^b
f(x)\,dx
}
\]

when the limit exists and is finite.

%%%%%%%%%%%%%%%%%%%%%%%%%%%%%%%%%%%%%%%%%%%%%%%%%%%%%%%%%%%%
\subsection{Improper Integral at a Finite Endpoint}
\label{subsec:finite-endpoint-improper-integral}
%%%%%%%%%%%%%%%%%%%%%%%%%%%%%%%%%%%%%%%%%%%%%%%%%%%%%%%%%%%%

If \(f\) becomes unbounded as \(x\to a^+\),

\[
\boxed{
\int_a^b f(x)\,dx
=
\lim_{c\to a^+}
\int_c^b
f(x)\,dx.
}
\]

%%%%%%%%%%%%%%%%%%%%%%%%%%%%%%%%%%%%%%%%%%%%%%%%%%%%%%%%%%%%
\subsection{\(p\)-Integral on \([1,\infty)\)}
\label{subsec:p-integral-infinity-reference}
%%%%%%%%%%%%%%%%%%%%%%%%%%%%%%%%%%%%%%%%%%%%%%%%%%%%%%%%%%%%

The integral

\[
\int_1^\infty
\frac{
1
}{
x^p
}
\,dx
\]

converges exactly when

\[
\boxed{
p>1.
}
\]

Thus

\[
\boxed{
\int_1^\infty
\frac{
dx
}{
x^p
}
\begin{cases}
\text{converges}, & p>1,\\
\text{diverges}, & p\leq1.
\end{cases}
}
\]

%%%%%%%%%%%%%%%%%%%%%%%%%%%%%%%%%%%%%%%%%%%%%%%%%%%%%%%%%%%%
\subsection{\(p\)-Integral Near Zero}
\label{subsec:p-integral-zero-reference}
%%%%%%%%%%%%%%%%%%%%%%%%%%%%%%%%%%%%%%%%%%%%%%%%%%%%%%%%%%%%

The integral

\[
\int_0^1
\frac{
1
}{
x^p
}
\,dx
\]

converges exactly when

\[
\boxed{
p<1.
}
\]

%%%%%%%%%%%%%%%%%%%%%%%%%%%%%%%%%%%%%%%%%%%%%%%%%%%%%%%%%%%%
\subsection{Worked Example: Improper Integral}
\label{subsec:worked-improper-integral}
%%%%%%%%%%%%%%%%%%%%%%%%%%%%%%%%%%%%%%%%%%%%%%%%%%%%%%%%%%%%

\begin{workedexamplebox}{Determine Convergence}

Evaluate

\[
\int_1^\infty
\frac1{x^2}
\,dx.
\]

Write

\[
\int_1^\infty
\frac1{x^2}
\,dx
=
\lim_{b\to\infty}
\int_1^b
x^{-2}\,dx.
\]

Then

\[
\int_1^b
x^{-2}\,dx
=
\left[
-\frac1x
\right]_1^b
=
1-\frac1b.
\]

Therefore,

\begin{answerbox}
\[
\boxed{
\int_1^\infty
\frac1{x^2}
\,dx
=
1.
}
\]
\end{answerbox}

\end{workedexamplebox}

%%%%%%%%%%%%%%%%%%%%%%%%%%%%%%%%%%%%%%%%%%%%%%%%%%%%%%%%%%%%
\subsection{Comparison Test for Improper Integrals}
\label{subsec:comparison-test-improper-integrals}
%%%%%%%%%%%%%%%%%%%%%%%%%%%%%%%%%%%%%%%%%%%%%%%%%%%%%%%%%%%%

Suppose

\[
0\leq f(x)\leq g(x)
\]

for sufficiently large \(x\).

Then:

\begin{itemize}
    \item if
    \[
    \int_a^\infty g(x)\,dx
    \]
    converges, then
    \[
    \int_a^\infty f(x)\,dx
    \]
    converges;

    \item if
    \[
    \int_a^\infty f(x)\,dx
    \]
    diverges, then
    \[
    \int_a^\infty g(x)\,dx
    \]
    diverges.
\end{itemize}

%%%%%%%%%%%%%%%%%%%%%%%%%%%%%%%%%%%%%%%%%%%%%%%%%%%%%%%%%%%%
\subsection{Limit Comparison for Improper Integrals}
\label{subsec:limit-comparison-integrals}
%%%%%%%%%%%%%%%%%%%%%%%%%%%%%%%%%%%%%%%%%%%%%%%%%%%%%%%%%%%%

For positive functions \(f,g\), if

\[
\boxed{
\lim_{x\to\infty}
\frac{
f(x)
}{
g(x)
}
=
L
}
\]

with

\[
0<L<\infty,
\]

then

\[
\int_a^\infty f(x)\,dx
\]

and

\[
\int_a^\infty g(x)\,dx
\]

have the same convergence behavior.

%%%%%%%%%%%%%%%%%%%%%%%%%%%%%%%%%%%%%%%%%%%%%%%%%%%%%%%%%%%%
\subsection{Gaussian Integral}
\label{subsec:gaussian-integral-reference}
%%%%%%%%%%%%%%%%%%%%%%%%%%%%%%%%%%%%%%%%%%%%%%%%%%%%%%%%%%%%

One of the most important definite integrals is

\[
\boxed{
\int_{-\infty}^{\infty}
e^{-x^2}
\,dx
=
\sqrt{\pi}.
}
\]

Consequently,

\[
\boxed{
\int_0^\infty
e^{-x^2}
\,dx
=
\frac{
\sqrt{\pi}
}{
2
}.
}
\]

No elementary antiderivative exists for \(e^{-x^2}\).

%%%%%%%%%%%%%%%%%%%%%%%%%%%%%%%%%%%%%%%%%%%%%%%%%%%%%%%%%%%%
\subsection{Scaled Gaussian Integral}
\label{subsec:scaled-gaussian-integral}
%%%%%%%%%%%%%%%%%%%%%%%%%%%%%%%%%%%%%%%%%%%%%%%%%%%%%%%%%%%%

For \(a>0\),

\[
\boxed{
\int_{-\infty}^{\infty}
e^{-ax^2}
\,dx
=
\sqrt{
\frac{\pi}{a}
}.
}
\]

Thus

\[
\boxed{
\int_0^\infty
e^{-ax^2}
\,dx
=
\frac12
\sqrt{
\frac{\pi}{a}
}.
}
\]

%%%%%%%%%%%%%%%%%%%%%%%%%%%%%%%%%%%%%%%%%%%%%%%%%%%%%%%%%%%%
\subsection{Normal Distribution Normalization}
\label{subsec:normal-distribution-integral}
%%%%%%%%%%%%%%%%%%%%%%%%%%%%%%%%%%%%%%%%%%%%%%%%%%%%%%%%%%%%

The standard normal density satisfies

\[
\boxed{
\int_{-\infty}^{\infty}
\frac{
1
}{
\sqrt{2\pi}
}
e^{-x^2/2}
\,dx
=
1.
}
\]

More generally,

\[
\boxed{
\int_{-\infty}^{\infty}
\frac{
1
}{
\sigma\sqrt{2\pi}
}
e^{-(x-\mu)^2/(2\sigma^2)}
\,dx
=
1
}
\]

for \(\sigma>0\).

%%%%%%%%%%%%%%%%%%%%%%%%%%%%%%%%%%%%%%%%%%%%%%%%%%%%%%%%%%%%
\subsection{Gamma Function}
\label{subsec:gamma-function-integral-reference}
%%%%%%%%%%%%%%%%%%%%%%%%%%%%%%%%%%%%%%%%%%%%%%%%%%%%%%%%%%%%

For \(s>0\),

\[
\boxed{
\Gamma(s)
=
\int_0^\infty
x^{s-1}e^{-x}
\,dx.
}
\]

For positive integers \(n\),

\[
\boxed{
\Gamma(n)
=
(n-1)!.
}
\]

%%%%%%%%%%%%%%%%%%%%%%%%%%%%%%%%%%%%%%%%%%%%%%%%%%%%%%%%%%%%
\subsection{Gamma Recurrence}
\label{subsec:gamma-recurrence-integral}
%%%%%%%%%%%%%%%%%%%%%%%%%%%%%%%%%%%%%%%%%%%%%%%%%%%%%%%%%%%%

Integration by parts gives

\[
\boxed{
\Gamma(s+1)
=
s\Gamma(s).
}
\]

In particular,

\[
\boxed{
\Gamma(n+1)
=
n!.
}
\]

%%%%%%%%%%%%%%%%%%%%%%%%%%%%%%%%%%%%%%%%%%%%%%%%%%%%%%%%%%%%
\subsection{Gamma at One-Half}
\label{subsec:gamma-one-half}
%%%%%%%%%%%%%%%%%%%%%%%%%%%%%%%%%%%%%%%%%%%%%%%%%%%%%%%%%%%%

A famous identity is

\[
\boxed{
\Gamma
\left(
\frac12
\right)
=
\sqrt{\pi}.
}
\]

This connects the Gamma function with the Gaussian integral.

%%%%%%%%%%%%%%%%%%%%%%%%%%%%%%%%%%%%%%%%%%%%%%%%%%%%%%%%%%%%
\subsection{Beta Function}
\label{subsec:beta-function-integral-reference}
%%%%%%%%%%%%%%%%%%%%%%%%%%%%%%%%%%%%%%%%%%%%%%%%%%%%%%%%%%%%

For \(p,q>0\),

\[
\boxed{
B(p,q)
=
\int_0^1
t^{p-1}
(1-t)^{q-1}
\,dt.
}
\]

%%%%%%%%%%%%%%%%%%%%%%%%%%%%%%%%%%%%%%%%%%%%%%%%%%%%%%%%%%%%
\subsection{Beta-Gamma Relation}
\label{subsec:beta-gamma-relation}
%%%%%%%%%%%%%%%%%%%%%%%%%%%%%%%%%%%%%%%%%%%%%%%%%%%%%%%%%%%%

The Beta and Gamma functions satisfy

\[
\boxed{
B(p,q)
=
\frac{
\Gamma(p)\Gamma(q)
}{
\Gamma(p+q)
}.
}
\]

%%%%%%%%%%%%%%%%%%%%%%%%%%%%%%%%%%%%%%%%%%%%%%%%%%%%%%%%%%%%
\subsection{Sine Power Integral}
\label{subsec:sine-power-definite-integral}
%%%%%%%%%%%%%%%%%%%%%%%%%%%%%%%%%%%%%%%%%%%%%%%%%%%%%%%%%%%%

For suitable \(m>-1\),

\[
\boxed{
\int_0^\pi
\sin^m x
\,dx
=
\sqrt{\pi}
\,
\frac{
\Gamma
\left(
\frac{m+1}{2}
\right)
}{
\Gamma
\left(
\frac{m}{2}+1
\right)
}.
}
\]

%%%%%%%%%%%%%%%%%%%%%%%%%%%%%%%%%%%%%%%%%%%%%%%%%%%%%%%%%%%%
\subsection{Wallis-Type Integral}
\label{subsec:wallis-integral-reference}
%%%%%%%%%%%%%%%%%%%%%%%%%%%%%%%%%%%%%%%%%%%%%%%%%%%%%%%%%%%%

Define

\[
I_n
=
\int_0^{\pi/2}
\sin^n x
\,dx.
\]

Then

\[
\boxed{
I_n
=
\frac{
n-1
}{
n
}
I_{n-2}
}
\]

for \(n\geq2\).

The initial values are

\[
\boxed{
I_0=\frac{\pi}{2},
\qquad
I_1=1.
}
\]

%%%%%%%%%%%%%%%%%%%%%%%%%%%%%%%%%%%%%%%%%%%%%%%%%%%%%%%%%%%%
\subsection{Fourier Orthogonality Integrals}
\label{subsec:fourier-orthogonality-integrals}
%%%%%%%%%%%%%%%%%%%%%%%%%%%%%%%%%%%%%%%%%%%%%%%%%%%%%%%%%%%%

For positive integers \(m,n\),

\[
\boxed{
\int_{-\pi}^{\pi}
\sin(mx)\sin(nx)
\,dx
=
\begin{cases}
0, & m\neq n,\\
\pi, & m=n,
\end{cases}
}
\]

and

\[
\boxed{
\int_{-\pi}^{\pi}
\cos(mx)\cos(nx)
\,dx
=
\begin{cases}
0, & m\neq n,\\
\pi, & m=n.
\end{cases}
}
\]

Also,

\[
\boxed{
\int_{-\pi}^{\pi}
\sin(mx)\cos(nx)
\,dx
=
0.
}
\]

%%%%%%%%%%%%%%%%%%%%%%%%%%%%%%%%%%%%%%%%%%%%%%%%%%%%%%%%%%%%
\subsection{Dirichlet-Type Sine Integral}
\label{subsec:dirichlet-sine-integral}
%%%%%%%%%%%%%%%%%%%%%%%%%%%%%%%%%%%%%%%%%%%%%%%%%%%%%%%%%%%%

An important improper integral is

\[
\boxed{
\int_0^\infty
\frac{
\sin x
}{
x
}
\,dx
=
\frac{\pi}{2}.
}
\]

This integral converges conditionally rather than absolutely.

%%%%%%%%%%%%%%%%%%%%%%%%%%%%%%%%%%%%%%%%%%%%%%%%%%%%%%%%%%%%
\subsection{Laplace-Type Integral}
\label{subsec:laplace-type-integral}
%%%%%%%%%%%%%%%%%%%%%%%%%%%%%%%%%%%%%%%%%%%%%%%%%%%%%%%%%%%%

For \(a>0\),

\[
\boxed{
\int_0^\infty
e^{-ax}
\,dx
=
\frac1a.
}
\]

More generally,

\[
\boxed{
\int_0^\infty
x^n e^{-ax}
\,dx
=
\frac{
n!
}{
a^{n+1}
}
}
\]

for nonnegative integer \(n\).

%%%%%%%%%%%%%%%%%%%%%%%%%%%%%%%%%%%%%%%%%%%%%%%%%%%%%%%%%%%%
\subsection{Moment-Type Gamma Integral}
\label{subsec:moment-gamma-integral}
%%%%%%%%%%%%%%%%%%%%%%%%%%%%%%%%%%%%%%%%%%%%%%%%%%%%%%%%%%%%

For \(a>0\) and \(s>0\),

\[
\boxed{
\int_0^\infty
x^{s-1}e^{-ax}
\,dx
=
\frac{
\Gamma(s)
}{
a^s
}.
}
\]

%%%%%%%%%%%%%%%%%%%%%%%%%%%%%%%%%%%%%%%%%%%%%%%%%%%%%%%%%%%%
\subsection{Common Definite Integrals}
\label{subsec:common-definite-integrals}
%%%%%%%%%%%%%%%%%%%%%%%%%%%%%%%%%%%%%%%%%%%%%%%%%%%%%%%%%%%%

\[
\boxed{
\int_0^1 x^n\,dx
=
\frac{1}{n+1},
\qquad
n>-1
}
\]

\[
\boxed{
\int_0^\pi
\sin x\,dx
=
2
}
\]

\[
\boxed{
\int_0^{2\pi}
\sin x\,dx
=
0
}
\]

\[
\boxed{
\int_0^{2\pi}
\cos x\,dx
=
0
}
\]

\[
\boxed{
\int_0^{\pi/2}
\sin x\,dx
=
1
}
\]

\[
\boxed{
\int_0^{\pi/2}
\cos x\,dx
=
1
}
\]

\[
\boxed{
\int_0^{\pi/2}
\sin^2x\,dx
=
\frac{\pi}{4}
}
\]

\[
\boxed{
\int_0^{\pi/2}
\cos^2x\,dx
=
\frac{\pi}{4}.
}
\]

%%%%%%%%%%%%%%%%%%%%%%%%%%%%%%%%%%%%%%%%%%%%%%%%%%%%%%%%%%%%
\subsection{Integral of a Derivative}
\label{subsec:integral-of-derivative}
%%%%%%%%%%%%%%%%%%%%%%%%%%%%%%%%%%%%%%%%%%%%%%%%%%%%%%%%%%%%

If \(F\) is differentiable and \(F'\) is integrable,

\[
\boxed{
\int_a^b
F'(x)
\,dx
=
F(b)-F(a).
}
\]

This expresses total change as accumulated instantaneous change.

%%%%%%%%%%%%%%%%%%%%%%%%%%%%%%%%%%%%%%%%%%%%%%%%%%%%%%%%%%%%
\subsection{Net Change Formula}
\label{subsec:net-change-integral}
%%%%%%%%%%%%%%%%%%%%%%%%%%%%%%%%%%%%%%%%%%%%%%%%%%%%%%%%%%%%

If

\[
Q'(t)=r(t),
\]

then

\[
\boxed{
Q(b)-Q(a)
=
\int_a^b r(t)\,dt.
}
\]

This applies to:

\begin{itemize}
    \item displacement from velocity;
    \item velocity change from acceleration;
    \item population change from growth rate;
    \item accumulated cost from marginal cost;
    \item mass from density;
    \item probability from density.
\end{itemize}

%%%%%%%%%%%%%%%%%%%%%%%%%%%%%%%%%%%%%%%%%%%%%%%%%%%%%%%%%%%%
\subsection{Probability Density Integrals}
\label{subsec:probability-density-integrals}
%%%%%%%%%%%%%%%%%%%%%%%%%%%%%%%%%%%%%%%%%%%%%%%%%%%%%%%%%%%%

For a continuous probability density \(f\),

\[
\boxed{
f(x)\geq0
}
\]

and

\[
\boxed{
\int_{-\infty}^{\infty}
f(x)\,dx
=
1.
}
\]

Probability over an interval is

\[
\boxed{
\Pr(a\leq X\leq b)
=
\int_a^b
f(x)\,dx.
}
\]

%%%%%%%%%%%%%%%%%%%%%%%%%%%%%%%%%%%%%%%%%%%%%%%%%%%%%%%%%%%%
\subsection{Expectation as an Integral}
\label{subsec:expectation-integral-reference}
%%%%%%%%%%%%%%%%%%%%%%%%%%%%%%%%%%%%%%%%%%%%%%%%%%%%%%%%%%%%

For continuous random variable \(X\) with density \(f\),

\[
\boxed{
\mathbb{E}[X]
=
\int_{-\infty}^{\infty}
x f(x)\,dx
}
\]

when the integral exists.

More generally,

\[
\boxed{
\mathbb{E}[g(X)]
=
\int_{-\infty}^{\infty}
g(x)f(x)\,dx.
}
\]

%%%%%%%%%%%%%%%%%%%%%%%%%%%%%%%%%%%%%%%%%%%%%%%%%%%%%%%%%%%%
\subsection{Mass From Density}
\label{subsec:mass-density-integral-reference}
%%%%%%%%%%%%%%%%%%%%%%%%%%%%%%%%%%%%%%%%%%%%%%%%%%%%%%%%%%%%

For linear density \(\rho(x)\),

\[
\boxed{
M
=
\int_a^b
\rho(x)\,dx.
}
\]

For volumetric density \(\rho(\mathbf{x})\),

\[
\boxed{
M
=
\int_{\Omega}
\rho
\,dV.
}
\]

%%%%%%%%%%%%%%%%%%%%%%%%%%%%%%%%%%%%%%%%%%%%%%%%%%%%%%%%%%%%
\subsection{Center of Mass in One Dimension}
\label{subsec:center-mass-one-dimensional}
%%%%%%%%%%%%%%%%%%%%%%%%%%%%%%%%%%%%%%%%%%%%%%%%%%%%%%%%%%%%

For density \(\rho(x)\),

\[
\boxed{
M
=
\int_a^b
\rho(x)\,dx,
}
\]

and center of mass is

\[
\boxed{
\bar{x}
=
\frac{
1
}{
M
}
\int_a^b
x\rho(x)\,dx.
}
\]

%%%%%%%%%%%%%%%%%%%%%%%%%%%%%%%%%%%%%%%%%%%%%%%%%%%%%%%%%%%%
\subsection{Work as an Integral}
\label{subsec:work-integral-reference}
%%%%%%%%%%%%%%%%%%%%%%%%%%%%%%%%%%%%%%%%%%%%%%%%%%%%%%%%%%%%

For force \(F(x)\) acting along a line,

\[
\boxed{
W
=
\int_a^b
F(x)\,dx.
}
\]

In vector form along a curve \(C\),

\[
\boxed{
W
=
\int_C
\mathbf{F}\cdot d\mathbf{r}.
}
\]

%%%%%%%%%%%%%%%%%%%%%%%%%%%%%%%%%%%%%%%%%%%%%%%%%%%%%%%%%%%%
\subsection{Line Integrals}
\label{subsec:line-integrals-reference}
%%%%%%%%%%%%%%%%%%%%%%%%%%%%%%%%%%%%%%%%%%%%%%%%%%%%%%%%%%%%

For scalar field \(f\) along parameterized curve

\[
\mathbf{r}(t),
\qquad
a\leq t\leq b,
\]

\[
\boxed{
\int_C
f\,ds
=
\int_a^b
f(\mathbf{r}(t))
\|\mathbf{r}'(t)\|
\,dt.
}
\]

%%%%%%%%%%%%%%%%%%%%%%%%%%%%%%%%%%%%%%%%%%%%%%%%%%%%%%%%%%%%
\subsection{Vector Line Integral}
\label{subsec:vector-line-integral-reference}
%%%%%%%%%%%%%%%%%%%%%%%%%%%%%%%%%%%%%%%%%%%%%%%%%%%%%%%%%%%%

For vector field \(\mathbf{F}\),

\[
\boxed{
\int_C
\mathbf{F}\cdot d\mathbf{r}
=
\int_a^b
\mathbf{F}(\mathbf{r}(t))
\cdot
\mathbf{r}'(t)
\,dt.
}
\]

%%%%%%%%%%%%%%%%%%%%%%%%%%%%%%%%%%%%%%%%%%%%%%%%%%%%%%%%%%%%
\subsection{Double Integrals}
\label{subsec:double-integral-reference}
%%%%%%%%%%%%%%%%%%%%%%%%%%%%%%%%%%%%%%%%%%%%%%%%%%%%%%%%%%%%

For a region \(R\subseteq\mathbb{R}^2\),

\[
\boxed{
\iint_R
f(x,y)
\,dA
}
\]

represents two-dimensional accumulation.

For a type-I region,

\[
R
=
\{
(x,y):
a\leq x\leq b,
g_1(x)\leq y\leq g_2(x)
\},
\]

\[
\boxed{
\iint_R
f\,dA
=
\int_a^b
\int_{g_1(x)}^{g_2(x)}
f(x,y)
\,dy\,dx.
}
\]

%%%%%%%%%%%%%%%%%%%%%%%%%%%%%%%%%%%%%%%%%%%%%%%%%%%%%%%%%%%%
\subsection{Triple Integrals}
\label{subsec:triple-integral-reference}
%%%%%%%%%%%%%%%%%%%%%%%%%%%%%%%%%%%%%%%%%%%%%%%%%%%%%%%%%%%%

For a three-dimensional region \(\Omega\),

\[
\boxed{
\iiint_\Omega
f(x,y,z)
\,dV
}
\]

represents volumetric accumulation.

For \(f=1\),

\[
\boxed{
V
=
\iiint_\Omega
1\,dV.
}
\]

%%%%%%%%%%%%%%%%%%%%%%%%%%%%%%%%%%%%%%%%%%%%%%%%%%%%%%%%%%%%
\subsection{Polar-Area Element}
\label{subsec:polar-area-element-reference}
%%%%%%%%%%%%%%%%%%%%%%%%%%%%%%%%%%%%%%%%%%%%%%%%%%%%%%%%%%%%

Under

\[
x=r\cos\theta,
\qquad
y=r\sin\theta,
\]

the area element is

\[
\boxed{
dA
=
r\,dr\,d\theta.
}
\]

Thus

\[
\boxed{
\iint_R
f(x,y)\,dA
=
\iint
f(r\cos\theta,r\sin\theta)
r\,dr\,d\theta.
}
\]

%%%%%%%%%%%%%%%%%%%%%%%%%%%%%%%%%%%%%%%%%%%%%%%%%%%%%%%%%%%%
\subsection{Cylindrical Volume Element}
\label{subsec:cylindrical-volume-element}
%%%%%%%%%%%%%%%%%%%%%%%%%%%%%%%%%%%%%%%%%%%%%%%%%%%%%%%%%%%%

In cylindrical coordinates,

\[
\boxed{
dV
=
r\,dr\,d\theta\,dz.
}
\]

%%%%%%%%%%%%%%%%%%%%%%%%%%%%%%%%%%%%%%%%%%%%%%%%%%%%%%%%%%%%
\subsection{Spherical Volume Element}
\label{subsec:spherical-volume-element}
%%%%%%%%%%%%%%%%%%%%%%%%%%%%%%%%%%%%%%%%%%%%%%%%%%%%%%%%%%%%

Using spherical coordinates

\[
x
=
\rho\sin\phi\cos\theta,
\]

\[
y
=
\rho\sin\phi\sin\theta,
\]

\[
z
=
\rho\cos\phi,
\]

the volume element is

\[
\boxed{
dV
=
\rho^2\sin\phi
\,d\rho\,d\phi\,d\theta.
}
\]

%%%%%%%%%%%%%%%%%%%%%%%%%%%%%%%%%%%%%%%%%%%%%%%%%%%%%%%%%%%%
\subsection{Jacobian Change of Variables}
\label{subsec:jacobian-change-variables-integral}
%%%%%%%%%%%%%%%%%%%%%%%%%%%%%%%%%%%%%%%%%%%%%%%%%%%%%%%%%%%%

For a coordinate transformation

\[
(x,y)
=
T(u,v),
\]

\[
\boxed{
dx\,dy
=
\left|
\frac{
\partial(x,y)
}{
\partial(u,v)
}
\right|
du\,dv.
}
\]

The absolute Jacobian determinant accounts for local area scaling.

%%%%%%%%%%%%%%%%%%%%%%%%%%%%%%%%%%%%%%%%%%%%%%%%%%%%%%%%%%%%
\subsection{Quick Core Antiderivative Table}
\label{subsec:quick-core-antiderivative-table}
%%%%%%%%%%%%%%%%%%%%%%%%%%%%%%%%%%%%%%%%%%%%%%%%%%%%%%%%%%%%

\[
\boxed{
\int x^n\,dx
=
\frac{x^{n+1}}{n+1}+C,
\qquad
n\neq-1
}
\]

\[
\boxed{
\int\frac1x\,dx
=
\ln|x|+C
}
\]

\[
\boxed{
\int e^x\,dx
=
e^x+C
}
\]

\[
\boxed{
\int a^x\,dx
=
\frac{a^x}{\ln a}+C
}
\]

\[
\boxed{
\int\ln x\,dx
=
x\ln x-x+C
}
\]

\[
\boxed{
\int\frac{dx}{1+x^2}
=
\arctan x+C
}
\]

\[
\boxed{
\int\frac{dx}{\sqrt{1-x^2}}
=
\arcsin x+C.
}
\]

%%%%%%%%%%%%%%%%%%%%%%%%%%%%%%%%%%%%%%%%%%%%%%%%%%%%%%%%%%%%
\subsection{Quick Composite Pattern Table}
\label{subsec:quick-composite-integral-patterns}
%%%%%%%%%%%%%%%%%%%%%%%%%%%%%%%%%%%%%%%%%%%%%%%%%%%%%%%%%%%%

\[
\boxed{
\int
u^n u'\,dx
=
\frac{
u^{n+1}
}{
n+1
}
+C,
\qquad
n\neq-1
}
\]

\[
\boxed{
\int
\frac{u'}{u}
\,dx
=
\ln|u|+C
}
\]

\[
\boxed{
\int
e^u u'\,dx
=
e^u+C
}
\]

\[
\boxed{
\int
\cos u\,u'\,dx
=
\sin u+C
}
\]

\[
\boxed{
\int
\sin u\,u'\,dx
=
-\cos u+C
}
\]

\[
\boxed{
\int
\frac{
u'
}{
1+u^2
}
\,dx
=
\arctan u+C
}
\]

\[
\boxed{
\int
\frac{
u'
}{
\sqrt{1-u^2}
}
\,dx
=
\arcsin u+C.
}
\]

%%%%%%%%%%%%%%%%%%%%%%%%%%%%%%%%%%%%%%%%%%%%%%%%%%%%%%%%%%%%
\subsection{Common Errors}
\label{subsec:integral-tables-common-errors}
%%%%%%%%%%%%%%%%%%%%%%%%%%%%%%%%%%%%%%%%%%%%%%%%%%%%%%%%%%%%

\subsubsection{Forgetting the Constant of Integration}

For an indefinite integral,

\[
\boxed{
\int f(x)\,dx
=
F(x)+C.
}
\]

The \(+C\) represents the entire family of antiderivatives.

%%%%%%%%%%%%%%%%%%%%%%%%%%%%%%%%%%%%%%%%%%%%%%%%%%%%%%%%%%%%
\subsubsection{Adding \(+C\) to a Definite Integral}

A definite integral evaluates to a number.

Write

\[
\boxed{
\int_a^b
f(x)\,dx
=
F(b)-F(a),
}
\]

not

\[
F(b)-F(a)+C.
\]

%%%%%%%%%%%%%%%%%%%%%%%%%%%%%%%%%%%%%%%%%%%%%%%%%%%%%%%%%%%%
\subsubsection{Using the Power Rule When \(n=-1\)}

The formula

\[
\frac{
x^{n+1}
}{
n+1
}
\]

does not apply when

\[
n=-1.
\]

Instead,

\[
\boxed{
\int x^{-1}\,dx
=
\ln|x|+C.
}
\]

%%%%%%%%%%%%%%%%%%%%%%%%%%%%%%%%%%%%%%%%%%%%%%%%%%%%%%%%%%%%
\subsubsection{Forgetting the Chain-Rule Factor in Reverse}

For example,

\[
\int
\cos(3x)\,dx
\]

is not

\[
\sin(3x)+C.
\]

Because

\[
\frac{d}{dx}
\sin(3x)
=
3\cos(3x),
\]

the correct result is

\[
\boxed{
\int
\cos(3x)\,dx
=
\frac13
\sin(3x)+C.
}
\]

%%%%%%%%%%%%%%%%%%%%%%%%%%%%%%%%%%%%%%%%%%%%%%%%%%%%%%%%%%%%
\subsubsection{Incorrect Integral of a Product}

In general,

\[
\boxed{
\int
f(x)g(x)\,dx
\neq
\left(
\int f(x)\,dx
\right)
\left(
\int g(x)\,dx
\right).
}
\]

Products require other methods such as substitution or integration by parts.

%%%%%%%%%%%%%%%%%%%%%%%%%%%%%%%%%%%%%%%%%%%%%%%%%%%%%%%%%%%%
\subsubsection{Incorrect Integral of a Quotient}

In general,

\[
\boxed{
\int
\frac{
f(x)
}{
g(x)
}
\,dx
\neq
\frac{
\int f(x)\,dx
}{
\int g(x)\,dx
}.
}
\]

%%%%%%%%%%%%%%%%%%%%%%%%%%%%%%%%%%%%%%%%%%%%%%%%%%%%%%%%%%%%
\subsubsection{Forgetting Absolute Value in Logarithmic Antiderivatives}

The correct real antiderivative is

\[
\boxed{
\int
\frac1x
\,dx
=
\ln|x|+C.
}
\]

Likewise,

\[
\boxed{
\int
\frac{
u'
}{
u
}
\,dx
=
\ln|u|+C.
}
\]

%%%%%%%%%%%%%%%%%%%%%%%%%%%%%%%%%%%%%%%%%%%%%%%%%%%%%%%%%%%%
\subsubsection{Failing to Change Bounds During Substitution}

If the integral is definite and one remains in the \(u\)-variable, the bounds
must also be transformed.

Do not insert \(x\)-bounds into an antiderivative expressed in \(u\).

%%%%%%%%%%%%%%%%%%%%%%%%%%%%%%%%%%%%%%%%%%%%%%%%%%%%%%%%%%%%
\subsubsection{Forgetting the Differential in Substitution}

If

\[
u=g(x),
\]

then

\[
\boxed{
du=g'(x)\,dx.
}
\]

A valid substitution must account for the differential factor.

%%%%%%%%%%%%%%%%%%%%%%%%%%%%%%%%%%%%%%%%%%%%%%%%%%%%%%%%%%%%
\subsubsection{Assuming Every Integral Has an Elementary Antiderivative}

For example,

\[
\boxed{
\int e^{-x^2}\,dx
}
\]

cannot be expressed using only elementary functions.

Special functions or numerical methods may be required.

%%%%%%%%%%%%%%%%%%%%%%%%%%%%%%%%%%%%%%%%%%%%%%%%%%%%%%%%%%%%
\subsubsection{Treating Improper Integrals as Ordinary Integrals}

An integral such as

\[
\int_1^\infty f(x)\,dx
\]

must be interpreted using a limit.

One cannot substitute \(\infty\) into an antiderivative as though infinity
were a real number.

%%%%%%%%%%%%%%%%%%%%%%%%%%%%%%%%%%%%%%%%%%%%%%%%%%%%%%%%%%%%
\subsubsection{Assuming Positive Integrands Always Give Finite Area}

For example,

\[
\frac1x
\]

is positive on \([1,\infty)\), but

\[
\boxed{
\int_1^\infty
\frac1x
\,dx
}
\]

diverges.

%%%%%%%%%%%%%%%%%%%%%%%%%%%%%%%%%%%%%%%%%%%%%%%%%%%%%%%%%%%%
\subsubsection{Confusing Signed Area With Geometric Area}

A definite integral gives signed area.

If \(f(x)<0\), then

\[
\int f(x)\,dx
\]

contributes negatively.

Geometric area requires integrating absolute vertical distance.

%%%%%%%%%%%%%%%%%%%%%%%%%%%%%%%%%%%%%%%%%%%%%%%%%%%%%%%%%%%%
\subsubsection{Using Trigonometric Substitution Without Domain Care}

When converting back from

\[
x=a\sin\theta
\]

or

\[
x=a\sec\theta,
\]

triangle relationships and sign restrictions must be consistent with the
chosen interval for \(\theta\).

%%%%%%%%%%%%%%%%%%%%%%%%%%%%%%%%%%%%%%%%%%%%%%%%%%%%%%%%%%%%
\subsubsection{Forgetting Polynomial Division}

Before partial fractions, if

\[
\deg P
\geq
\deg Q,
\]

perform polynomial division.

%%%%%%%%%%%%%%%%%%%%%%%%%%%%%%%%%%%%%%%%%%%%%%%%%%%%%%%%%%%%
\subsubsection{Using Partial Fractions Before Factoring Completely}

The denominator should be factored over the relevant number system before the
partial-fraction form is chosen.

%%%%%%%%%%%%%%%%%%%%%%%%%%%%%%%%%%%%%%%%%%%%%%%%%%%%%%%%%%%%
\subsubsection{Confusing Antiderivatives With Areas}

An indefinite integral is a family of functions.

A definite integral is a number representing net accumulation.

They are related but not identical concepts.

%%%%%%%%%%%%%%%%%%%%%%%%%%%%%%%%%%%%%%%%%%%%%%%%%%%%%%%%%%%%
\subsection{Integration Strategy}
\label{subsec:integration-strategy-reference}
%%%%%%%%%%%%%%%%%%%%%%%%%%%%%%%%%%%%%%%%%%%%%%%%%%%%%%%%%%%%

A practical integration workflow is:

\begin{enumerate}
    \item simplify algebraically when useful;
    \item check for a basic antiderivative pattern;
    \item look for an inner function and its derivative;
    \item try substitution;
    \item identify products suited for integration by parts;
    \item use trigonometric identities for trig powers or products;
    \item use trigonometric substitution for quadratic radicals;
    \item factor rational-function denominators;
    \item use polynomial division when necessary;
    \item use partial fractions for proper rational functions;
    \item complete the square for quadratic expressions;
    \item consider special functions for nonelementary antiderivatives;
    \item for definite integrals, exploit symmetry or periodicity;
    \item for improper integrals, replace the problematic endpoint by a limit;
    \item differentiate the final indefinite result to verify it.
\end{enumerate}

%%%%%%%%%%%%%%%%%%%%%%%%%%%%%%%%%%%%%%%%%%%%%%%%%%%%%%%%%%%%
\subsection{Exercises}
\label{subsec:integral-table-exercises}
%%%%%%%%%%%%%%%%%%%%%%%%%%%%%%%%%%%%%%%%%%%%%%%%%%%%%%%%%%%%

\begin{exercise}[1]
Evaluate

\[
\int x^5\,dx.
\]
\end{exercise}

\begin{exercise}[1]
Evaluate

\[
\int
(4x^3-2x+7)
\,dx.
\]
\end{exercise}

\begin{exercise}[1]
Evaluate

\[
\int\frac1x\,dx.
\]
\end{exercise}

\begin{exercise}[1]
Evaluate

\[
\int e^x\,dx.
\]
\end{exercise}

\begin{exercise}[1]
Evaluate

\[
\int 3^x\,dx.
\]
\end{exercise}

\begin{exercise}[1]
Evaluate

\[
\int\sin x\,dx.
\]
\end{exercise}

\begin{exercise}[1]
Evaluate

\[
\int\cos x\,dx.
\]
\end{exercise}

\begin{exercise}[1]
Evaluate

\[
\int\sec^2x\,dx.
\]
\end{exercise}

\begin{exercise}[1]
Evaluate

\[
\int\csc^2x\,dx.
\]
\end{exercise}

\begin{exercise}[2]
Evaluate

\[
\int
(2x+1)^5
\,dx.
\]
\end{exercise}

\begin{exercise}[2]
Evaluate

\[
\int
2x e^{x^2}
\,dx.
\]
\end{exercise}

\begin{exercise}[2]
Evaluate

\[
\int
\frac{
2x
}{
x^2+5
}
\,dx.
\]
\end{exercise}

\begin{exercise}[2]
Evaluate

\[
\int
\frac{
3
}{
1+9x^2
}
\,dx.
\]
\end{exercise}

\begin{exercise}[2]
Evaluate

\[
\int
\frac{
dx
}{
\sqrt{1-x^2}
}.
\]
\end{exercise}

\begin{exercise}[2]
Evaluate

\[
\int
\cos(5x)
\,dx.
\]
\end{exercise}

\begin{exercise}[2]
Evaluate

\[
\int
\sin(7x)
\,dx.
\]
\end{exercise}

\begin{exercise}[2]
Evaluate

\[
\int\tan x\,dx.
\]
\end{exercise}

\begin{exercise}[2]
Evaluate

\[
\int\cot x\,dx.
\]
\end{exercise}

\begin{exercise}[2]
Evaluate

\[
\int\sec x\,dx.
\]
\end{exercise}

\begin{exercise}[2]
Evaluate

\[
\int\csc x\,dx.
\]
\end{exercise}

\begin{exercise}[2]
Evaluate

\[
\int\sinh x\,dx.
\]
\end{exercise}

\begin{exercise}[2]
Evaluate

\[
\int\tanh x\,dx.
\]
\end{exercise}

\begin{exercise}[2]
Use substitution to evaluate

\[
\int
x(x^2+1)^4
\,dx.
\]
\end{exercise}

\begin{exercise}[2]
Use substitution to evaluate

\[
\int
e^{3x+2}
\,dx.
\]
\end{exercise}

\begin{exercise}[2]
Use substitution to evaluate

\[
\int
\frac{
x
}{
\sqrt{x^2+4}
}
\,dx.
\]
\end{exercise}

\begin{exercise}[2]
Evaluate

\[
\int_0^1
2x(x^2+1)^3
\,dx
\]

by changing the bounds.
\end{exercise}

\begin{exercise}[2]
Use integration by parts to evaluate

\[
\int xe^x\,dx.
\]
\end{exercise}

\begin{exercise}[2]
Use integration by parts to evaluate

\[
\int x\sin x\,dx.
\]
\end{exercise}

\begin{exercise}[2]
Use integration by parts to evaluate

\[
\int\ln x\,dx.
\]
\end{exercise}

\begin{exercise}[3]
Evaluate

\[
\int x^2e^x\,dx.
\]
\end{exercise}

\begin{exercise}[3]
Evaluate

\[
\int e^x\sin x\,dx.
\]
\end{exercise}

\begin{exercise}[3]
Evaluate

\[
\int e^{2x}\cos3x\,dx.
\]
\end{exercise}

\begin{exercise}[2]
Evaluate

\[
\int\sin^3x\cos x\,dx.
\]
\end{exercise}

\begin{exercise}[2]
Evaluate

\[
\int\sin x\cos^4x\,dx.
\]
\end{exercise}

\begin{exercise}[2]
Evaluate

\[
\int\sin^2x\,dx.
\]
\end{exercise}

\begin{exercise}[2]
Evaluate

\[
\int\cos^2x\,dx.
\]
\end{exercise}

\begin{exercise}[3]
Evaluate

\[
\int\sin^4x\,dx.
\]
\end{exercise}

\begin{exercise}[3]
Evaluate

\[
\int\tan^3x\sec^2x\,dx.
\]
\end{exercise}

\begin{exercise}[3]
Evaluate

\[
\int\tan x\sec^3x\,dx.
\]
\end{exercise}

\begin{exercise}[2]
Use product-to-sum to evaluate

\[
\int
\sin4x\cos2x
\,dx.
\]
\end{exercise}

\begin{exercise}[3]
Evaluate using trigonometric substitution:

\[
\int
\sqrt{9-x^2}
\,dx.
\]
\end{exercise}

\begin{exercise}[3]
Evaluate

\[
\int
\frac{
dx
}{
\sqrt{x^2+4}
}.
\]
\end{exercise}

\begin{exercise}[3]
Evaluate

\[
\int
\frac{
dx
}{
x^2+6x+13
}.
\]
\end{exercise}

\begin{exercise}[2]
Use partial fractions to evaluate

\[
\int
\frac{
1
}{
x(x+1)
}
\,dx.
\]
\end{exercise}

\begin{exercise}[3]
Evaluate

\[
\int
\frac{
3x+5
}{
x^2+3x+2
}
\,dx.
\]
\end{exercise}

\begin{exercise}[3]
Use polynomial division and partial fractions to evaluate a rational function
whose numerator degree is at least the denominator degree.
\end{exercise}

\begin{exercise}[2]
Evaluate

\[
\int_0^1
x^3\,dx.
\]
\end{exercise}

\begin{exercise}[2]
Evaluate

\[
\int_0^\pi
\sin x\,dx.
\]
\end{exercise}

\begin{exercise}[2]
Evaluate

\[
\int_{-\pi}^{\pi}
x^3\cos x
\,dx
\]

using symmetry.
\end{exercise}

\begin{exercise}[2]
Evaluate

\[
\int_{-2}^{2}
(x^4+1)
\,dx
\]

using even symmetry.
\end{exercise}

\begin{exercise}[2]
Find the average value of

\[
f(x)=x^2
\]

on \([0,3]\).
\end{exercise}

\begin{exercise}[2]
Find the area between

\[
y=x
\]

and

\[
y=x^2
\]

on \([0,1]\).
\end{exercise}

\begin{exercise}[2]
Use disks to find the volume obtained by rotating

\[
y=x
\]

on \([0,1]\) about the \(x\)-axis.
\end{exercise}

\begin{exercise}[2]
Use cylindrical shells to compute a simple solid-of-revolution volume.
\end{exercise}

\begin{exercise}[3]
Find the arc length of

\[
y=x^{3/2}
\]

over a specified interval.
\end{exercise}

\begin{exercise}[2]
Evaluate

\[
\int_1^\infty
\frac1{x^2}
\,dx.
\]
\end{exercise}

\begin{exercise}[2]
Determine whether

\[
\int_1^\infty
\frac1x
\,dx
\]

converges.
\end{exercise}

\begin{exercise}[2]
Determine convergence of

\[
\int_1^\infty
\frac{
1
}{
x^{3/2}
}
\,dx.
\]
\end{exercise}

\begin{exercise}[2]
Determine convergence of

\[
\int_0^1
\frac{
1
}{
\sqrt{x}
}
\,dx.
\]
\end{exercise}

\begin{exercise}[3]
Use comparison to test an improper integral.
\end{exercise}

\begin{exercise}[3]
Use limit comparison to test an improper integral.
\end{exercise}

\begin{exercise}[3]
Verify

\[
\Gamma(1)=1.
\]
\end{exercise}

\begin{exercise}[3]
Use integration by parts to prove

\[
\Gamma(s+1)
=
s\Gamma(s).
\]
\end{exercise}

\begin{exercise}[3]
Show that

\[
\Gamma(n+1)=n!
\]

for positive integers \(n\).
\end{exercise}

\begin{exercise}[3]
Evaluate

\[
\int_0^\infty
xe^{-x}
\,dx
\]

using the Gamma function.
\end{exercise}

\begin{exercise}[3]
Evaluate

\[
\int_0^\infty
x^2e^{-2x}
\,dx.
\]
\end{exercise}

\begin{exercise}[3]
Verify that the standard normal density integrates to \(1\) using the
Gaussian integral.
\end{exercise}

\begin{exercise}[3]
Use the Beta-Gamma relation to evaluate a simple Beta integral.
\end{exercise}

\begin{exercise}[3]
Use the Wallis recurrence to evaluate

\[
\int_0^{\pi/2}
\sin^4x
\,dx.
\]
\end{exercise}

\begin{exercise}[3]
Show that

\[
\int_{-\pi}^{\pi}
\sin(mx)\cos(nx)
\,dx
=
0
\]

for positive integers \(m,n\).
\end{exercise}

\begin{exercise}[3]
Find the mass of a rod with density

\[
\rho(x)
=
1+x
\]

on \([0,2]\).
\end{exercise}

\begin{exercise}[3]
Find the center of mass of the same rod.
\end{exercise}

\begin{exercise}[3]
Compute work done by a variable force

\[
F(x)=kx
\]

from \(x=0\) to \(x=a\).
\end{exercise}

\begin{exercise}[3]
For a continuous density \(f\), express

\[
\Pr(a\leq X\leq b)
\]

as an integral.
\end{exercise}

\begin{exercise}[3]
Compute

\[
\mathbb{E}[X]
\]

for a simple continuous density.
\end{exercise}

\begin{exercise}[3]
Evaluate a double integral over a rectangular region.
\end{exercise}

\begin{exercise}[3]
Evaluate a double integral in polar coordinates.
\end{exercise}

\begin{exercise}[3]
Use cylindrical coordinates to compute the volume of a cylinder.
\end{exercise}

\begin{exercise}[3]
Use spherical coordinates to derive the volume of a sphere.
\end{exercise}

\begin{exercise}[4]
Derive integration by parts from the product rule.
\end{exercise}

\begin{exercise}[4]
Derive the substitution rule from the chain rule.
\end{exercise}

\begin{exercise}[4]
Derive the integral of secant.
\end{exercise}

\begin{exercise}[4]
Derive the integral of cosecant.
\end{exercise}

\begin{exercise}[4]
Derive the reduction formula for

\[
\int\sin^n x\,dx.
\]
\end{exercise}

\begin{exercise}[4]
Derive the reduction formula for

\[
\int\cos^n x\,dx.
\]
\end{exercise}

\begin{exercise}[4]
Derive the formula

\[
\int
e^{ax}\sin bx
\,dx.
\]
\end{exercise}

\begin{exercise}[4]
Derive the formula

\[
\int
e^{ax}\cos bx
\,dx.
\]
\end{exercise}

\begin{exercise}[4]
Derive

\[
\int
\sqrt{a^2-x^2}
\,dx.
\]
\end{exercise}

\begin{exercise}[4]
Prove the \(p\)-integral convergence conditions.
\end{exercise}

\begin{exercise}[4]
Derive the Gaussian integral by squaring the integral and converting to polar
coordinates.
\end{exercise}

\begin{exercise}[4]
Derive

\[
\Gamma
\left(
\frac12
\right)
=
\sqrt{\pi}.
\]
\end{exercise}

\begin{exercise}[4]
Derive the Beta-Gamma relation.
\end{exercise}

\begin{exercise}[4]
Study why

\[
e^{-x^2}
\]

has no elementary antiderivative.
\end{exercise}

%%%%%%%%%%%%%%%%%%%%%%%%%%%%%%%%%%%%%%%%%%%%%%%%%%%%%%%%%%%%
\subsection{Conceptual Summary}
\label{subsec:integral-tables-summary}
%%%%%%%%%%%%%%%%%%%%%%%%%%%%%%%%%%%%%%%%%%%%%%%%%%%%%%%%%%%%

An indefinite integral represents a family of antiderivatives:

\[
\boxed{
\int f(x)\,dx
=
F(x)+C.
}
\]

The fundamental power rule is

\[
\boxed{
\int x^n\,dx
=
\frac{
x^{n+1}
}{
n+1
}
+
C,
\qquad
n\neq-1.
}
\]

The exceptional logarithmic case is

\[
\boxed{
\int\frac1x\,dx
=
\ln|x|+C.
}
\]

Important exponential and trigonometric formulas include

\[
\boxed{
\int e^x\,dx
=
e^x+C,
}
\]

\[
\boxed{
\int\sin x\,dx
=
-\cos x+C,
}
\]

and

\[
\boxed{
\int\cos x\,dx
=
\sin x+C.
}
\]

Substitution reverses the chain rule:

\[
\boxed{
\int
f(g(x))g'(x)
\,dx
=
\int f(u)\,du.
}
\]

Integration by parts reverses the product rule:

\[
\boxed{
\int u\,dv
=
uv-\int v\,du.
}
\]

Definite integrals satisfy

\[
\boxed{
\int_a^b f(x)\,dx
=
F(b)-F(a)
}
\]

when \(F'=f\).

Improper integrals require limits.

Special integrals such as

\[
\boxed{
\int_{-\infty}^{\infty}
e^{-x^2}
\,dx
=
\sqrt{\pi}
}
\]

lead naturally to special functions, probability distributions, and advanced
analysis.

The central practical principle is

\[
\boxed{
\text{identify the structural pattern before choosing an integration method}.
}
\]

%%%%%%%%%%%%%%%%%%%%%%%%%%%%%%%%%%%%%%%%%%%%%%%%%%%%%%%%%%%%
\subsection{Section Connections}
\label{subsec:integral-tables-connections}
%%%%%%%%%%%%%%%%%%%%%%%%%%%%%%%%%%%%%%%%%%%%%%%%%%%%%%%%%%%%

\chapterconnection{
Integral Tables consolidates antiderivatives and definite-integral patterns
used throughout Calculus, Analysis, Differential Equations, Probability,
Statistics, Geometry, Physics, Continuum Mechanics, Fourier Analysis,
Optimization, and Mathematical Modeling. Substitution reverses the Chain
Rule from Section~\ref{sec:derivative-tables}; integration by parts reverses
the Product Rule; trigonometric identities from
Section~\ref{sec:trigonometric-identities} support trigonometric integration;
improper integrals connect with limit laws; Gaussian, Gamma, and Beta
integrals connect with Probability and Mathematical Statistics; and
multivariable integration connects with Geometry and Mathematical Physics.
The next reference section, Series Expansions, collects geometric, binomial,
Taylor, Maclaurin, exponential, logarithmic, trigonometric, hyperbolic,
Fourier-related, and asymptotic expansions.
}

%%%%%%%%%%%%%%%%%%%%%%%%%%%%%%%%%%%%%%%%%%%%%%%%%%%%%%%%%%%%
% SECTION SOURCE TRACKING
%%%%%%%%%%%%%%%%%%%%%%%%%%%%%%%%%%%%%%%%%%%%%%%%%%%%%%%%%%%%

%------------------------------------------------------------
% 14.10 Integral Tables
%------------------------------------------------------------
%
% Source basis:
%   - Standard single-variable and multivariable calculus.
%   - Standard trigonometric, exponential, logarithmic,
%     rational, improper, Gamma, Beta, Gaussian,
%     probability, and geometric integral formulas.
%
% Used for:
%   - integral notation;
%   - indefinite integrals;
%   - definite integrals;
%   - Fundamental Theorem of Calculus;
%   - constant and linearity rules;
%   - power rule;
%   - logarithmic antiderivative;
%   - exponential integrals;
%   - trigonometric integrals;
%   - inverse-trigonometric forms;
%   - hyperbolic forms;
%   - substitution;
%   - definite substitution;
%   - integration by parts;
%   - cyclic integration by parts;
%   - trigonometric power integrals;
%   - power reduction;
%   - product-to-sum integration;
%   - trigonometric substitution;
%   - completing the square;
%   - partial fractions;
%   - Weierstrass substitution;
%   - reduction formulas;
%   - definite-integral identities;
%   - symmetry;
%   - periodic functions;
%   - average value;
%   - areas and volumes;
%   - arc length;
%   - improper integrals;
%   - p-integrals;
%   - comparison tests;
%   - Gaussian integral;
%   - normal-density normalization;
%   - Gamma function;
%   - Beta function;
%   - Fourier orthogonality;
%   - Laplace-type integrals;
%   - probability and expectation integrals;
%   - density and center-of-mass integrals;
%   - work integrals;
%   - line, double, and triple integrals;
%   - coordinate Jacobians;
%   - common errors;
%   - applications and exercises.
%
% Notes:
%   - Series Expansions is developed next in Section 14.11.
%   - Transform Tables follows in Section 14.12.
%
%%%%%%%%%%%%%%%%%%%%%%%%%%%%%%%%%%%%%%%%%%%%%%%%%%%%%%%%%%%%